\documentclass[11pt]{amsart}
\usepackage{amsmath,amssymb,amsthm}
\usepackage{booktabs}
\usepackage[margin=1.15in]{geometry}
\usepackage[colorlinks=true,linkcolor=blue,citecolor=blue,urlcolor=blue]{hyperref}
\usepackage{microtype}
\emergencystretch=1.5em
\raggedbottom

\newtheorem{theorem}{Theorem}[section]
\newtheorem{proposition}[theorem]{Proposition}
\newtheorem{corollary}[theorem]{Corollary}
\newtheorem{lemma}[theorem]{Lemma}
\newtheorem{conjecture}[theorem]{Conjecture}
\newtheorem{problem}[theorem]{Problem}
\newtheorem{question}[theorem]{Question}
\theoremstyle{definition}
\newtheorem{definition}[theorem]{Definition}
\newtheorem{example}[theorem]{Example}
\theoremstyle{remark}
\newtheorem{remark}[theorem]{Remark}
\newtheorem{observation}[theorem]{Observation}

\newcommand{\PP}{\mathbb{P}}
\newcommand{\CC}{\mathbb{C}}
\newcommand{\ZZ}{\mathbb{Z}}
\newcommand{\QQ}{\mathbb{Q}}
\newcommand{\RR}{\mathbb{R}}
\newcommand{\NN}{\mathbb{N}}
\newcommand{\gfrak}{\mathfrak{g}}
\newcommand{\Dscr}{\mathfrak{D}}
\newcommand{\chat}{\hat{c}}
\newcommand{\Bal}{\operatorname{Balanced}}
\newcommand{\Crit}{\operatorname{Crit}}
\newcommand{\Mbar}{\overline{\mathcal{M}}}
\newcommand{\Ascr}{\mathcal{A}}
\newcommand{\Om}{\Omega_{r,s}}

\hyphenation{scat-ter-ing tropi-cal tropi-cal-izes singu-lar-ity de-scen-dent}

\title[Boundary-ray factorization in open FJRW theory]{Boundary-ray
factorization and geometric wall-crossing\\ in open FJRW theory of
$x^{r}+y^{s}$}

\author{Bernd Johannes Wuebben}
\subjclass[2020]{14N35 (primary); 14J33, 14T90, 53D45 (secondary)}
\keywords{Open FJRW theory, wall-crossing, tropical vertex,
Landau--Ginzburg models, boundary conditions}
\date{Initial version: August 24, 2026. This version: October 9, 2026}
\hypersetup{
  pdftitle={Boundary-ray factorization and geometric wall-crossing in open FJRW theory of x\string^r+y\string^s},
  pdfauthor={Bernd Johannes Wuebben},
  pdfsubject={Boundary-ray factorization and geometric wall-crossing in open FJRW theory of x\string^r+y\string^s},
  pdfkeywords={open FJRW theory, Landau--Ginzburg models, wall-crossing, tropical vertex, boundary conditions}
}

\begin{document}

\begin{abstract}
We prove that in the genus-zero open FJRW theory of $x^r+y^s$ of Gross, Kelly
and Tessler, the wall term of each critical disk lies on the ray spanned by
its numbers of boundary marked points, and that the wall-crossing
transformation between the two extremal systems of invariants, in which on each set of invariants related by walls only the coefficient of least, respectively greatest, $x$-degree can be nonzero, factors uniquely into automorphisms
supported on these rays and ordered by slope.  The proof rests on the
Kontsevich--Soibelman ordered-product lemma, applied with the filtration by
the number of markings since the ray assignment is not a grading, and on the
simply transitive action of the wall-crossing group once a torus direction
is added to its Lie algebra.  With descendents the factorization holds in
every finite truncation, where it is realized by a homotopy of boundary
conditions whose jumps occur ray by ray; in general no single homotopy in
the normal form of the Gross--Kelly--Tessler extension lemma does this.  In
each finite truncation the invariants between consecutive rays are obtained
by transport across the rays from constants analogous to theta-function
coefficients, not all of which wall-crossing produces.
\end{abstract}

\maketitle
\enlargethispage{3pt}

\section{Introduction}
\label{sec:intro}

\subsection{The theory}
\label{ssec:theory}
Gross, Kelly and Tessler \cite{GKT} (henceforth GKT) defined genus-zero open
FJRW invariants
\[
\nu_{\Gamma,\mathbf d}
=\Big\langle\prod_{i}\tau_{d_{i}}^{(a_{i},b_{i})}\,
\sigma_{1}^{\beta_{1}}\sigma_{2}^{\beta_{2}}\sigma_{12}\Big\rangle^{\mathfrak s,o}
\]
for the Landau--Ginzburg model $(W=x^{r}+y^{s},\ \mu_{r}\times\mu_{s})$: relative
Euler numbers of descendent Witten bundles over moduli of $W$-spin disks, which are
real orbifolds with corners. Here $\tau_{d}^{(a,b)}$ is an internal marking of
twist $(a,b)$, $0\le a\le r-2$, $0\le b\le s-2$, with descendent exponent $d$, and
$\sigma_1$, $\sigma_2$, $\sigma_{12}$ are boundary markings of twists $(r{-}2,0)$,
$(0,s{-}2)$ and $(r{-}2,s{-}2)$, the last being the \emph{fully twisted} one.  The
invariants depend on a family $\mathfrak s$
of canonical multisections (boundary conditions), unlike the open $r$-spin
intersection numbers of Buryak, Clader and Tessler, whose construction GKT
generalize and which do not depend on this choice \cite[Thm.~3.17]{BCT21}.

GKT record the combinatorial type of a disk by a \emph{graph} $\Gamma$.  The
graphs that carry invariants are the \emph{balanced} ones: they have
one $\sigma_{12}$-point, and the descendent Witten bundle has rank equal to the
dimension
of their moduli space, so that its relative Euler number is a number once
boundary conditions are fixed \cite[Def.~2.32]{GKT}.  A \emph{critical}
graph has no $\sigma_{12}$-point and rank one more than its dimension
\cite[Def.~2.33]{GKT}, so that a homotopy of boundary conditions transverse to
the zero section vanishes at finitely many points of $[0,1]$ times its moduli
space.

On the boundary of a balanced moduli space the homotopy can vanish only
on the strata where a critical graph and a balanced graph meet at a boundary
node, and the invariants change only through these zeros.  GKT assemble the
change at each time at which such zeros occur into an automorphism of the
potential below: the exponential of the wall terms of the critical graphs, the
Hamiltonian vector fields attached to them (see the conventions in
\S\ref{ssec:results}), weighted, up to explicit signs, by their zero counts.
The total change is the time-ordered product of these exponentials
\cite[proofs of Thms.~5.4 and~5.9]{GKT}; this jump formula is restated in
\S\ref{sec:presentation}.

The invariants are collected in the
potential \cite[(4.11)]{GKT}
\[
W^{\nu}=\sum_{J,\mathbf d}\ \sum_{\Gamma}
\frac{(-1)^{|J|-1}}{|\mathrm{Aut}(J,\mathbf d)|}\,
\nu_{\Gamma,\mathbf d}\,u_{J,\mathbf d}\,x^{\beta_1(\Gamma)}y^{\beta_2(\Gamma)}
\;=\;x^{r}+y^{s}+\cdots,
\]
the sums over multisets $J$ of internal twists $(a_j,b_j)$ with descendent
vectors $\mathbf d$ and over the balanced graphs $\Gamma$ with these
internal markings.  Here
$u_{J,\mathbf d}=\prod_{j}t_{a_{j},b_{j},d_{j}}$ is a monomial in deformation
parameters $t_{a,b,d}$ dual to the insertions $\tau_{d}^{(a,b)}$, and
$\beta(\Gamma)=(\beta_1(\Gamma),\beta_2(\Gamma))$ is the \emph{boundary degree} of
$\Gamma$, its numbers of boundary markings of types $\sigma_1$ and $\sigma_2$; the
empty multiset contributes $x^r+y^s$, by the normalization $\nu=-1$ for
balanced graphs without internal markings, and over the labeled ring of
\cite[Def.~4.6]{GKT}
the same formula holds without the symmetry factor.  Up to the sign and
symmetry factor, which are constant along each diagonal (see the conventions in
\S\ref{ssec:results}),
the coefficient of $u_{J,\mathbf d}x^{\beta_1}y^{\beta_2}$ is thus the invariant
of the balanced graph of boundary degree $\beta$ over $(J,\mathbf d)$.

Three of GKT's results organize the dependence on boundary conditions
completely. First, the systems
of invariants arising from canonical families form a set $\Om(I)$, for $I$ a
finite set of internal marking labels as in \cite{GKT}, characterized by
explicit polynomial conditions, the \emph{chamber-index} conditions of
\cite[Def.~3.47, Cor.~5.14]{GKT}: normalizations for at most one internal
marking; the condition $\Ascr(J,\mathbf d,\nu)=(-1)^{d_1}$ when $J$ is a single
marking with descendent exponent $d_1$; and the condition
$\Ascr(J,\mathbf d,\nu)=0$ when $|J|\ge2$ and $N(J,\mathbf d)\ge0$, that is,
when balanced graphs over $(J,\mathbf d)$ exist (the integer
$N(J,\mathbf d)$ is defined in the conventions of \S\ref{ssec:results}, and
\S\ref{sec:algebra}
shows that this is the inequality of \cite[Def.~3.47]{GKT}).  Here $\Ascr$ is a
$\Gamma$-function-weighted polynomial in the invariants, recalled in
\eqref{eq:Apoly}.  A system satisfying these
conditions is a \emph{chamber index} in GKT's terminology; the word does not
refer to a chamber of the wall complement in the quadrant used below.  Since
these are exactly the systems in $\Om(I)$, we call them \emph{realizable}.

Second,
GKT state that their wall-crossing group acts faithfully and transitively on
these systems \cite[Thm.~4.27 and Cor.~5.14]{GKT}; transitivity requires the
marginal torus summands adjoined in \S\ref{ssec:results}
(\cite[Prop.~1.4 and Rem.~3.1]{note}, Proposition~\ref{prop:nontransitive}).
Third, the
potential $W^{\nu}$ of any system has oscillatory integrals over the good
basis, the cycles $\Xi_{(a,b)}$, $0\le a\le r-2$, $0\le b\le s-2$, with
$\int_{\Xi_{(a,b)}}x^{a'}y^{b'}e^{(x^r+y^s)/\hbar}\,dx\wedge dy
=\delta_{aa'}\delta_{bb'}$ for $(a',b')$ in the same range \cite[(4.3)]{GKT}, equal to the closed extended FJRW
generating function
\cite[Thm.~0.3]{GKT}, an open mirror theorem building on the closed
Landau--Ginzburg mirror symmetry of \cite{LLSS,HLSW}.

GKT remark \cite[Rem.~0.8]{GKT} that their wall-crossing group is
close in spirit to the tropical vertex group of Kontsevich--Soibelman \cite{KS06},
as used for Fano mirror symmetry in \cite{GPS10,Gross10} and in the log
Calabi--Yau setting \cite{GHK15}, and their survey raises
the problem of a tropical formulation of the open theory \cite{GKTsurvey}. Such a
formulation is not immediate. A Landau--Ginzburg model carries no target space and
hence no SYZ fibration, and Maher notes that without an SYZ fibration the
Gross--Siebert machinery does not apply directly to these models
\cite[Rem.~5.1.1]{Maher}.

This paper organizes the wall-crossing group of GKT by rays in the
quadrant of boundary degrees.  It proves that in each finite truncation each element of the
group factors uniquely as a slope-ordered product of factors supported on
single rays, the \emph{ray factors}, and computes the first factors; the
construction is algebraic and does not presuppose a degenerating target.  Here
the \emph{boundary ray} of a wall term is the ray spanned by the boundary
degree of its critical graph, that is, by its numbers of boundary marked
points; it is not one of the two edge rays of the quadrant, which we call
the \emph{edges} of the quadrant.  The paper then shows that in each finite
truncation the factorization of the \emph{transition element}
$g^{\min\to\max}$, the wall-crossing transformation between the two extremal
systems, in which the coefficients on every wall-bearing diagonal
(\S\ref{ssec:results}) vanish except at one end, is realized by homotopies of
boundary conditions; the ray factors, but not the chosen homotopies, are
compatible in the coefficientwise limit.  By \emph{geometric wall-crossing} we
mean this realization of the factorization by homotopies of boundary
conditions whose jumps occur ray by ray.

The paper also shows that the system of
open invariants at every slope cut, a direction in the quadrant taken on one
side of any wall it meets (\S\ref{ssec:results}), is obtained by transport
through the ray factors from initial coefficients, constants attached to the
monomials of the potential that do not depend on the cut.  The quadrant organizes
wall-crossing from one of its edges to the other; it does not by itself carry
the singular integral-affine structure of a Gross--Siebert base.  For
$x^4+y^4$ the companion paper \cite{annular} constructs a singular
integral-affine annulus and on it a wall structure, consistent at its two joints, whose singleton walls are the ray factors of $g^{\min\to\max}$ for a single first-descendent marking of twist $(1,0)$ or $(0,1)$, but whose first wall carrying one marking of each of these twists is not (Remark~\ref{rem:interpretation}).

\subsection{Results}
\label{ssec:results}
In order of importance, the main new results are the following.
\begin{enumerate}
\item[(a)] The wall term of a critical graph fixes exactly the monomials whose
exponents are proportional to its boundary degree, as a wall of a scattering
diagram fixes the monomials along it; its boundary ray is therefore the unique
ray whose monomials it fixes, and the walls are placed there
(Proposition~\ref{prop:placement}).  Along these rays the transition element
$g^{\min\to\max}$ between the two extremal systems factors without choices
(Theorem~\ref{thm:canonical}).  By contrast, the jump formula of \cite[proof of Thm.~5.9]{GKT} orders the total jump of a normal-form homotopy by homotopy time, with one
factor for each subset of internal markings and factors that mix boundary
rays; the times and their order are auxiliary choices of the extension lemma
(Remark~\ref{rem:timeorder}).  The factorization itself is the
Kontsevich--Soibelman ordered-product lemma, applied with the filtration by
the number of markings because the ray assignment is not a grading
(Theorem~\ref{thm:presentation}(ii)).
\item[(b)] In each finite coefficient quotient this factorization is realized
by concatenated homotopies of boundary conditions, one ray at a time
(Theorem~\ref{thm:georealization}), whereas a single homotopy in the normal
form of the GKT extension lemma makes all critical graphs with the same internal
markings jump at once and so cannot in general separate their rays
(Proposition~\ref{prop:obstruction} and Corollary~\ref{cor:nosimple}).
\item[(c)] The invariants at every slope cut are transported through the ray
factors from initial coefficients (Theorem~\ref{thm:transport}); nonzero
initial coefficients occur beyond those of $x^r$ and $y^s$, so the open theory
is strictly richer than the part generated from $x^r+y^s$ by wall-crossing
(Observation~\ref{obs:richer}).
\end{enumerate}
Theorems~\ref{thm:canonical}, \ref{thm:georealization} and~\ref{thm:transport}
below state the main parts of (a)--(c) in this order; they are preceded by
Theorem~\ref{thm:presentation}, whose part (ii) is the factorization lemma
used in (a).  Under \emph{Further results}, Theorem~\ref{thm:effective} makes
explicit on each diagonal how the conditions of GKT determine the invariants,
and Theorem~\ref{thm:trichotomy} classifies the wall terms by weight.

Two further results are recalled from the companion note \cite{note}, where
they are proved over the symmetric coefficient ring.  First, the transitivity
asserted in \cite[Thm.~4.27]{GKT} requires adjoining to the Lie algebra of
\cite[Def.~4.23]{GKT} the torus direction $x\partial_x-y\partial_y$ of the
wall terms of critical graphs with one boundary marking of each type, and the
enlarged group acts freely and transitively \cite[Prop.~1.4 and
Rem.~3.1]{note}; the proof extends the invariance of the polynomials $\Ascr$,
which GKT derive from a period formula stated for residues at most
$(r-2,s-2)$, to all residues through a Laplace integral.  Here the torsor property is proved over the labeled ring of \cite[Def.~4.6]{GKT} as well (Lemma~\ref{lem:enlarged}), as the realization theorem requires, using the all-residue invariance (Lemma~\ref{lem:allresidue}), whose Laplace computation \cite[proof of Prop.~1.4]{note} already carries out over that ring; and
Proposition~\ref{prop:nontransitive} adds an explicit primary instance of the
failure, for $r=s=4$ and two labels of twist $(2,2)$; \cite[Ex.~3.2]{note}
gives a descendent one.  Second, Theorem~\ref{thm:trichotomy}, including the
threshold (iii), is also in the companion note \cite[\S2, Lem.~4.2 and
Prop.~7.1]{note}.

The organizing fact is a dictionary between the wall-crossing algebra of
\cite{GKT} and the tropical vertex.  The wall term of a critical graph is the
monomial of its internal markings times the Hamiltonian vector field, for the
area form $dx\wedge dy$, of the boundary monomial $x^{\#\sigma_1}y^{\#\sigma_2}$
of the critical graph.  Since $dx\wedge dy=xy\,d\log x\wedge d\log y$, this
vector field is $(xy)^{-1}$ times the tropical-vertex vector field whose
exponent is the boundary degree
$(\#\sigma_1,\#\sigma_2)$ (Proposition~\ref{prop:shift}).  The shift
by $(1,1)$ between this exponent and the Lie bidegree, the exponent shift that
the wall term effects on monomials, is the removal of the
single fully twisted boundary point $\sigma_{12}$, which balanced graphs carry
and critical graphs do not.  Consequently the wall term of a critical graph lies
on its boundary ray, and all results below are organized by these rays.

\subsection*{Conventions}
Throughout, $\gfrak^{r,s}$ denotes the Landau--Ginzburg wall-crossing Lie algebra
of \cite[Def.~4.23]{GKT}: it is spanned by the \emph{wall terms}
$u_{J,\mathbf d}X_{k}$, $k=(k_1,k_2)$, with
\[
X_{k}=X_{k_{1},k_{2}}=x^{k_{1}}y^{k_{2}}\bigl((k_{2}{+}1)x\partial_{x}-(k_{1}{+}1)
y\partial_{y}\bigr)
\]
the Hamiltonian vector field of $x^{k_{1}+1}y^{k_{2}+1}$ for $dx\wedge dy$ (so
that $X_{k}=(xy)^{-1}g_{k+(1,1)}$ for the vector field
$g_\alpha=z^\alpha(\alpha_2x\partial_x-\alpha_1y\partial_y)$ of the tropical
vertex), and
$u_{J,\mathbf d}$ the deformation monomial of \S\ref{ssec:theory} (the
descendent index of $t_{a,b,d}$ is suppressed when $d=0$; \emph{primary} means
that all $d_j=0$; the \emph{deformation ideal} $\mathfrak m$ is the ideal
generated by the $t_{a,b,d}$), subject to the
congruences $k_{1}\equiv r(J)\ (\mathrm{mod}\ r)$, $k_{2}\equiv s(J)\
(\mathrm{mod}\ s)$ and the dimension constraint
\begin{equation}
\label{eq:dimconstraint}
sk_{1}+rk_{2}=m(J,\mathbf d)-rs,
\qquad
m(J,\mathbf d)=rs+\sum_{j}\bigl(sa_{j}+rb_{j}+rs(d_{j}-1)\bigr).
\end{equation}

We call $k$ the \emph{Lie bidegree}, or bidegree, of $X_k$; a critical graph
of boundary degree $\beta=(\beta_1,\beta_2)$ carries wall terms of Lie bidegree
$k=\beta-(1,1)$ (Proposition~\ref{prop:shift}).  Throughout, $k$ denotes a Lie
bidegree and $\beta$ a boundary degree.  The \emph{slope} of a nonzero
boundary degree $\beta$, and of the ray it spans, is
$\beta_2/\beta_1\in[0,\infty]$; the $\sigma_1$-axis has slope $0$ and the
$\sigma_2$-axis slope $\infty$.  A \emph{slope-ordered product} of
automorphisms supported on rays composes them in the order of the slopes of
their rays; in increasing order the factor of least slope acts first, so that
for factors $\exp(\eta_i)$ on rays $\rho_i$ with
$\operatorname{slope}(\rho_1)<\cdots<\operatorname{slope}(\rho_L)$ the product
is $\exp(\eta_L)\circ\cdots\circ\exp(\eta_1)$.

The right-hand side of
\eqref{eq:dimconstraint} is the \emph{Euler weight} $w(J,\mathbf d)$ of the
monomial $u_{J,\mathbf d}$.  Write
\[
 \sum_{j\in J}a_j=r(J)+r\ell_1,\quad 0\le r(J)<r,\qquad
 \sum_{j\in J}b_j=s(J)+s\ell_2,\quad 0\le s(J)<s,
\]
and set
\[
 N(J,\mathbf d)=\ell_1+\ell_2-|J|+1+\sum_{j\in J}d_j.
\]
By \cite[Not.~2.40]{GKT}, when $N\ge0$ there are $N+1$ balanced
graphs, whose coefficients in the potential we call the \emph{balanced
coefficients}, and when $N\ge1$ there are $N$ critical
graphs; below the respective thresholds those sets are empty.  For fixed
$(J,\mathbf d)$ the $N+1$ balanced coefficients form the
\emph{$(J,\mathbf d)$-diagonal}, also called the \emph{diagonal} of
$u_{J,\mathbf d}$ (Definition~\ref{def:diagonal}): the points
$(r(J)+pr,\ s(J)+(N-p)s)$, $0\le p\le N$, spaced by $(r,-s)$ along a line in
the boundary-degree plane.  The wall term of each of the $N$ intervening
critical graphs changes $x^r+y^s$ by a \emph{wall vector} supported on two
adjacent points of the diagonal \eqref{eq:wallvector}.  The diagonal and its monomial are
\emph{wall-bearing} if $N\ge1$, and \emph{marginal} if $w(J,\mathbf d)=0$.

We include the wall terms of bidegree $(0,0)$ that occur when $w=0$.  Thus
\[
\widetilde\gfrak^{\,r,s}=\gfrak^{r,s}\oplus
\bigoplus_{\substack{\varnothing\ne J,\;w(J,\mathbf d)=0,\;N(J,\mathbf d)\ge1\\
r(J)=s(J)=0}}
\QQ\,u_{J,\mathbf d}(x\partial_x-y\partial_y),
\qquad
\widetilde G^{r,s}=\exp(\widetilde\gfrak^{\,r,s}).
\]
Remark~\ref{rem:torus} proves that this is a Lie algebra; its marginal torus
summands are omitted from \cite[Def.~4.23]{GKT}.  The enlarged group acts
freely and transitively on $\Om(I)$, the \emph{torsor action}
(Lemma~\ref{lem:enlarged}; \cite[Prop.~1.4]{note}), while the transitivity
asserted in \cite[Thm.~4.27]{GKT} fails without these summands as soon as a
marginal critical graph occurs (\cite[Rem.~3.1]{note};
Remark~\ref{rem:torus}, Proposition~\ref{prop:nontransitive}).  Any two
realizable systems $\nu,\nu'$ are thus related by a unique element of
$\widetilde G^{r,s}$, the \emph{transition element} from $\nu$ to $\nu'$.
All statements below are for the enlarged algebra and its group.

As in GKT, a \emph{family} is indexed by graphs, and a \emph{homotopy} is a
one-parameter deformation, of a single multisection or of a family.  A
\emph{family of simple canonical homotopies}
\cite[Defs.~3.36--3.37]{GKT}, which we also call a \emph{homotopy of boundary
conditions}, is a homotopy $(\mathfrak s_T)_{T\in[0,1]}$ of families of
canonical multisections, compatible on boundary strata.  The extension lemma
\cite[Lem.~3.40]{GKT}, recalled as Lemma~\ref{lem:gktextension}, constructs
such a homotopy between any two families of canonical multisections that are
transverse to the zero section.  The \emph{jumps} of such a homotopy are the
automorphisms by which its potential changes at the zeros of the homotopy on
critical boundary strata, given by the jump formula of \cite[proof of
Thm.~5.9]{GKT} restated in \S\ref{sec:presentation}; their time-ordered
product is its \emph{total jump}.

A
finite set $S$ of deformation monomials is \emph{divisor-closed} if every
monomial dividing an element of $S$ belongs to $S$; reduction to $S$ means
quotienting out all other deformation monomials.  The \emph{boundary ray} of
a wall term $u_{J,\mathbf d}X_{k}$ is $\RR_{\ge0}(k_1+1,k_2+1)$, the ray of the
boundary degree of its critical graph.  The \emph{boundary-degree quadrant}
$C_\partial$ is the closed first quadrant of the boundary-degree plane
(Definition~\ref{def:sector}), with edges the $\sigma_1$-axis and the
$\sigma_2$-axis.  A \emph{slope cut} is a direction in $C_\partial$, taken on
one side of the wall when a wall lies on that direction.  It plays the role of
the endpoint of a path in the path-ordered products of \cite{GPS10}, and,
because the descendent walls are dense in $C_\partial$
(Proposition~\ref{prop:localfinite}), slope cuts replace the chambers of a
finite wall structure (Definition~\ref{def:sector}).  The
\emph{$\sigma_1$-edge cut} and the \emph{$\sigma_2$-edge cut} are the cuts at
the two edges of $C_\partial$.

\subsection*{Main theorems}
The first theorem collects two general statements, on homotopies of boundary
conditions and on the group.

\begin{theorem}[Slope-ordered boundary-ray factorization; \S\ref{sec:presentation}]
\label{thm:presentation}
(i) (Path independence for boundary conditions.) For any homotopy of
boundary conditions supplied by \cite[Lem.~3.40]{GKT} with equal endpoints,
the time-ordered
product of its jumps, the critical-graph jump automorphisms of the jump formula
of \cite[proof of Thm.~5.9]{GKT} (restated in \S\ref{sec:presentation}), is
the identity; for such homotopies with arbitrary endpoints, the total jump
depends only on the endpoints. (ii) (Boundary-ray factorization.) In the primary
algebra, modulo every power of the deformation ideal, every element of the
enlarged wall-crossing group factors uniquely as a slope-ordered product of
exponentials supported on the boundary rays
$\RR_{\ge0}(k_{1}{+}1,k_{2}{+}1)$; the marginal torus terms belong to the
$(1,1)$-ray. The compatible factorizations define its completed boundary-ray
factorization. In the full descendent algebra the same factorization is taken
on finite divisor-closed coefficient sets and passed to the coefficientwise
limit.
\end{theorem}

Part (i) holds because the action is free (Lemma~\ref{lem:enlarged}): the
product of the jumps around a loop of boundary conditions fixes a potential,
hence is the identity.  It is a consistency statement in the space of
boundary conditions, distinct from consistency of the ray diagram in the
quadrant in the sense of \cite{GPS10}.  Part (ii) is the
Kontsevich--Soibelman ordered-product lemma \cite{KS06},
\cite[\S0.3 and Thm.~1.4]{GPS10}, here over the boundary rays, whose
directions are the Lie bidegrees shifted by $(1,1)$ and therefore do not form a
Lie grading.  The induction still works because each ray subalgebra is
abelian, the bracket being proportional to the determinant of the directions
of the two boundary rays, and the filtration by the number of markings
replaces the grading.

In each finite truncation the factorization (ii) of an element $g$ is a
finite wall structure of the shape of a scattering diagram, whose
path-ordered product from one edge of the quadrant to the other is $g$.  Apart
from the factors on the $(1,1)$-ray, its automorphisms are the rational-power
wall-crossings of Proposition~\ref{prop:shift}, which multiply $x$ and $y$ by
rational powers of binomials $1+c\,x^{k_1}y^{k_2}$ \cite[Lem.~4.21]{GKT},
not elements of the tropical vertex group of \cite{GPS10}.  For the transition
element $g^{\min\to\max}$ between the extremal systems of Theorem~\ref{thm:canonical}
below, this product is not the identity unless every ray factor is trivial.
Accordingly the diagram $\Dscr^{r,s}$ defined after Theorem~\ref{thm:canonical} is not
the consistent completion of an initial diagram: it is the
Kontsevich--Soibelman factorization of one prescribed element,
$g^{\min\to\max}$.  In the
graded setting of Kontsevich--Soibelman and Bridgeland \cite{KS06,Bridgeland},
ordered factorization over the rays of the grading identifies each group
element with the consistent scattering diagram whose path-ordered product
between two opposite chambers is that element; $\Dscr^{r,s}$ is the analogue
of that diagram for $g^{\min\to\max}$, read along a path from one edge of the
quadrant to the other, but the boundary rays do not come from a grading of
$\widetilde\gfrak^{\,r,s}$ (Proposition~\ref{prop:shift}), so the standard
bijection does not apply verbatim, and Theorem~\ref{thm:presentation}(ii)
supplies the factorization instead.

\begin{theorem}[Extremal systems and the boundary-ray factorization
$\Dscr^{r,s}$; \S\ref{sec:canonical}]
\label{thm:canonical}
On every wall-bearing diagonal, including a marginal diagonal, the wall action
on the $N{+}1$ balanced coefficients is triangular; consequently there exist
unique realizable systems $\nu^{\min}$ and $\nu^{\max}$ in which, on every such
diagonal, only the extreme coefficient survives. The factorization, in
increasing order of slope, of the unique group element $g^{\min\to\max}$
carrying $W^{\nu^{\min}}$ to $W^{\nu^{\max}}$ is a completed primary
boundary-ray factorization $\Dscr^{r,s}_{0}$, which involves no choices and is
computable order by order.
\end{theorem}

We call $\nu^{\min}$ and $\nu^{\max}$ the \emph{minimal} and \emph{maximal
systems}, together the \emph{extremal systems}: on each wall-bearing diagonal
the surviving coefficient is the one of least $x$-degree in $\nu^{\min}$ and
the one of greatest $x$-degree in $\nu^{\max}$.  For a finite labeled set $I$,
$\nu^{\min}_I$ and $\nu^{\max}_I$ denote the restrictions of $\nu^{\min}$ and
$\nu^{\max}$ to the labeled set $I$.  In the full descendent algebra the
proof gives compatible factorizations of $g^{\min\to\max}$ on all finite
divisor-closed coefficient sets (Theorem~\ref{thm:presentation}(ii)); their
coefficientwise limit, read on slope cuts, is the descendent extension
$\Dscr^{r,s}$ of $\Dscr^{r,s}_{0}$.  The wall action is triangular because
the $N$ wall vectors of a diagonal are supported on the $N$ pairs of adjacent
points with indices $\{p-1,p\}$, $1\le p\le N$, one pair each, with nonzero
entries at both points ($\lambda_p,\mu_p>0$ in \eqref{eq:wallvector}).  The existence
of $\nu^{\min}$ and $\nu^{\max}$ is the triangular elimination of \cite[proof of Thm.~4.27,
Step~3]{GKT}, with the torus direction adjoined; the new statement in
Theorem~\ref{thm:canonical} is the factorization of $g^{\min\to\max}$ along
the boundary rays, which involves no choices.  We write $\theta_{\rho}$ for
the factor of $\Dscr^{r,s}$ supported on the ray $\rho$ and call it the
\emph{ray factor} of $\rho$; it collects all critical-graph terms for which
$(k_{1}{+}1,k_{2}{+}1)$ lies on $\rho$ (Proposition~\ref{prop:shift}).  For
$x^{5}+y^{5}$ the walls of $\Dscr^{5,5}$ begin
\begin{gather*}
\theta_{(1,2)}=\exp\bigl(\tfrac1{25}\,t_{2,3}t_{3,3}X_{0,1}+\cdots\bigr),\qquad
\theta_{(2,1)}=\exp\bigl(\tfrac1{25}\,t_{3,2}t_{3,3}X_{1,0}+\cdots\bigr),\\
\theta_{(1,1)}=\exp\bigl(\tfrac1{50}\,t_{3,3}^{2}X_{1,1}
+\tfrac1{25}(t_{2,2}t_{3,3}+t_{2,3}t_{3,2})X_{0,0}+\cdots\bigr),
\end{gather*}
where the subscripts are the primitive boundary degrees on the rays.  The ray
factors are computed by solving $\nu^{\max}-\nu^{\min}$ diagonal by diagonal,
order by order, against the wall vectors (Example~\ref{ex:walls}).

The worked case $x^{5}+y^{5}$ (\S\ref{sec:canonical} and Appendix~\ref{sec:x5y5})
produces, beyond the ray factors above, the invariant relations and values
\begin{gather}
\langle(\tau_{0}^{(3,3)})^{2}\sigma_{1}\sigma_{2}^{6}\sigma_{12}\rangle
+\langle(\tau_{0}^{(3,3)})^{2}\sigma_{1}^{6}\sigma_{2}\sigma_{12}\rangle
=-\tfrac25,
\label{eq:firstwallrel}\\
\langle(\tau_{0}^{(3,3)})^{3}
\sigma_{1}^{4}\sigma_{2}^{4}\sigma_{12}\rangle=\tfrac15,\qquad
\langle\tau_{0}^{(1,2)}\tau_{0}^{(2,1)}\tau_{0}^{(2,2)}\sigma_{12}\rangle
=-\tfrac1{25},
\label{eq:newinvariants}
\end{gather}
the first a wall relation (only the sum is invariant; the difference changes
across the $(1,1)$-ray, which carries the first symmetric positive-weight primary wall
term $t_{3,3}^{2}X_{1,1}$), the second and third absolute: they live on
diagonals with no wall terms, the second on a broad diagonal, where
$r(J)=r-1$ or $s(J)=s-1$.  The \emph{two-insertion relation}, satisfied by every realizable system, takes the uniform form
\begin{equation}
\label{eq:seed}
\frac{\nu_{(r,0)}}{r}+\frac{\nu_{(0,s)}}{s}=-\frac1{rs}
\qquad\text{(e.g. }\langle\sigma_{1}^{r}\sigma_{12}\cdots\rangle
+\langle\sigma_{2}^{r}\sigma_{12}\cdots\rangle=-\tfrac1r\ \text{at }r=s),
\end{equation}
for every two-insertion diagonal $\{(a_{1},b_{1}),(a_{2},b_{2})\}$,
$|J|=2$, with $a_{1}+a_{2}=r$, $b_{1}+b_{2}=s$ (Remark~\ref{rem:ex55}). Here
$\nu_{\beta}$ denotes the invariant of boundary degree $\beta=(\beta_1,\beta_2)$,
with boundary insertions $\sigma_{1}^{\beta_{1}}\sigma_{2}^{\beta_{2}}\sigma_{12}$, and
the suppressed
insertions are the two interior markings $(a_1,b_1)$ and $(a_2,b_2)$.

The realization theorem uses the normal form of the extension lemma.  A
\emph{normal-form homotopy} is a homotopy of boundary conditions as
constructed in \cite[Lem.~3.40]{GKT} (Lemma~\ref{lem:gktextension}): each subset $A$
of the internal labels has one time $T_A$, and every jump of a critical graph
with internal marking set $A$ occurs at $T_A$ (\S\ref{sec:presentation}).  The
theorem also uses labeled sets large enough for symmetrization to be
injective.  In the universal finite-order statements below, the finite
labeled set $I$ is chosen \emph{saturated} at the order under consideration:
when one works modulo the $K$th power of the deformation ideal, it carries at
least $K-1$ labels of every primary twist, and, on a finite divisor-closed
set $S$, at least $n_{a,b}(S)$ labels of every twist $(a,b)$, with
$n_{a,b}(S)$ as in the next theorem.

\begin{theorem}[Realization by homotopies in a finite coefficient quotient;
  \S\ref{ssec:realization}]
\label{thm:georealization}
Let $S$ be a finite divisor-closed set of deformation monomials containing
$1$.  Choose a finite labeled set $I$ with at least
\[
n_{a,b}(S):=\max_{u\in S}\sum_{d\ge0}\operatorname{ord}_{t_{a,b,d}}(u)
\]
  labels of every twist $(a,b)$ occurring in $S$.  The reduction to $S$ of the factorization $\Dscr^{r,s}$ (Theorem~\ref{thm:canonical} and the paragraph after it), the \emph{quotient
factorization}, has finitely many rays $\rho_1,\dots,\rho_L$, listed in
increasing slope.  There is a homotopy of boundary conditions, that is, a
family of simple canonical homotopies in the sense of
\cite[Defs.~3.36--3.37]{GKT}, obtained by concatenating normal-form
homotopies, from a family of canonical multisections realizing
$\nu_I^{\min}$ to one realizing $\nu_I^{\max}$; every
time slice is a family of canonical multisections \cite[Def.~3.22]{GKT}.
After reduction to the quotient supported on $S$ every nonzero
jump is supported on a single boundary ray, the rays carrying
nonzero reduced jumps occur in increasing slope as the homotopy time
increases, and the total reduced jump on $\rho_{i}$ is the ray factor
$\theta_{\rho_{i}}$ of Theorem~\ref{thm:canonical}.
\end{theorem}

The proof (\S\ref{ssec:realization}) treats one ray at a time.  The
intermediate systems, obtained from $\nu_I^{\min}$ by applying the factors
$\theta_{\rho_1},\theta_{\rho_2},\dots$ in turn, are realized by canonical
families by \cite[Cor.~5.14]{GKT}, and consecutive ones are joined by a single
normal-form homotopy of \cite[Lem.~3.40]{GKT}.  By the jump formula of
\cite[proof of Thm.~5.9]{GKT} (\S\ref{sec:presentation}) the
total jump of such a homotopy is a time-ordered product of factors, one for
each subset of internal labels and each with exact label support on that
subset, and freeness of the action forces this total to equal the labeled
lift of $\theta_{\rho_i}$, its image under symmetrization in the labeled
coefficient ring of \cite[Def.~4.6]{GKT}.  A factorization with exact label supports is
unique, and this uniqueness, together with the abelianness of the subalgebra
of each ray, forces every internal jump of that homotopy onto the prescribed
ray; the concatenation may reuse a marking subset in later homotopies.  A final
normal-form homotopy joining the result to a realization of $\nu_I^{\max}$ is
trivial in the quotient.  The label count $n_{a,b}(S)$ makes symmetrization
injective \cite[Lem.~4.14]{GKT}, so the labeled statement descends to the
symmetric one.  The increasing order of slopes is the convention fixed in
Theorem~\ref{thm:canonical}: by Theorem~\ref{thm:presentation}(ii),
$g^{\min\to\max}$ also factors uniquely in decreasing order of slopes, and the
same construction realizes that factorization (Remark~\ref{rem:conjform}).

A concatenation is needed whenever a diagonal carries two wall terms with
nonzero reduced coefficients that receive no contribution from lower-order
terms, as for
$t_{3,3}^5$ in $x^5+y^5$ or for any descendent singleton $t_{a,b,2}$: a single
normal-form homotopy of \cite[Lem.~3.40]{GKT} cannot separate them, since all
wall terms of a diagonal jump at one time (Proposition~\ref{prop:obstruction}).
The exact first order at which two primary wall terms share a diagonal is
computed for every hyperbolic pair, that is, every pair with $rs>2r+2s$
(Proposition~\ref{prop:parabolicgeom},
Remark~\ref{rem:thresholdgeom}).  Realizing each boundary-ray factor by its
own normal-form homotopy avoids the obstruction.  On the punctured quadrant the
product of ray factors along a path transverse to the walls depends only on
the chambers of its angular endpoints
(Proposition~\ref{thm:positionobstruction7}).

The last main result is algebraic and concerns transport through the
quadrant $C_\partial$.  The \emph{transport}
$\Psi(q)$ to a slope cut $q$ is the slope-ordered product of the ray factors
from the $\sigma_1$-edge cut up to $q$, that is, the path-ordered product, in
the sense of \cite{GPS10}, along a path from the $\sigma_1$-edge to $q$, taken
coefficientwise (Definition~\ref{def:sector}).  The \emph{cut potential} at
$q$ is $W(q)=\Psi(q)(W^{\nu^{\min}})$.  Here $m$ denotes a monomial
$u_{J,\mathbf d}x^{\beta_1}y^{\beta_2}$ of the potential; the integer
$m(J,\mathbf d)$ of \eqref{eq:dimconstraint} always carries its arguments.
The \emph{anchor cut} $q_m$ of $m$ is the slope cut through the direction of
its exponent, on the side toward the $(1,1)$-direction
(Definition~\ref{def:anchor}), and the coefficient of $m$ at $q_m$ is its
\emph{anchored value}.

\begin{theorem}[The transport representation; \S\ref{sec:geometric}]
\label{thm:transport}
There are unique constants $\mathrm{init}_m$, one for each monomial $m$ of the
potential, such that every cut potential is their coefficientwise transport,
\[
W(q)=\sum_{m}\Psi(q)\,\Psi(q_{m})^{-1}(\mathrm{init}_{m}\,m),
\]
the sum over the monomials $m$ of the potential and $q_{m}$ the anchor cut
of $m$.  On every monomial $u_{J,\mathbf d}$ the difference between the
coefficient vector of $W(q)$ and the transported contributions of the
constants $\mathrm{init}_m$ of the proper divisor monomials is
\emph{independent of the cut} $q$, and it equals the vector of the constants
$\mathrm{init}_m$ on the diagonal of $u_{J,\mathbf d}$.
\end{theorem}

We call $\mathrm{init}_m$ the \emph{initial coefficient} of $m$; the initial
coefficients form the \emph{initial system}, and by the last statement they
are determined recursively in the divisibility order of deformation monomials.

The proof of Theorem~\ref{thm:transport} is an induction in the divisibility
order of deformation monomials: on
the diagonal of a monomial $u$, the change of the potential across a wall and
the change of the transported lower initial coefficients are the same universal
expression in lower-order data, so their difference does not depend on the
cut.  This part is formal.  What gives the representation content is that
the initial system is not generated by wall-crossing
(Observation~\ref{obs:richer}): transport through the ray factors from
$x^r+y^s$ alone does not produce the potential, because nonzero initial
coefficients occur beyond those of $x^r$ and $y^s$, for instance the invariant
$\tfrac15$ of \eqref{eq:newinvariants}, on a diagonal without wall terms, which
no ray action can produce.  In the smooth Fano case of \cite{Gross10,CPS}, by
contrast, the potential is generated by broken lines from its initial
monomials.  A coefficient whose initial
coefficient is zero is given at every cut by transport of lower initial
coefficients alone.  Two consequences, stated and proved after
Theorem~\ref{thm:transport7}, compute an initial coefficient directly when no
bending from another initial coefficient reaches the relevant cut.
Theorem~\ref{thm:transport} underlies the tropical description of the open
theory in \S\ref{sec:geometric}: the cut potentials are computed by bend chains
from the initial coefficients, where a \emph{bend chain} is a term of the
expansion of Lemma~\ref{lem:bendchain}, a sequence of wall terms applied to a
transported monomial along a monotone path in the slope line.  Bend chains
play the role of the broken lines of \cite{Gross10,CPS}, but there is no
affine manifold and no asymptotic direction.  In the language of theta
functions \cite{GHK15,GHKK}, written $\vartheta$ as there (the letter
$\theta_\rho$ for ray factors follows \cite{GPS10}), $\mathrm{init}_m$ plays
the role of the
coefficient of $\vartheta_m$ in an expansion $\sum_m\mathrm{init}_m\vartheta_m$
of the potential, which is therefore not the analogue of
$\vartheta_{(r,0)}+\vartheta_{(0,s)}$, and the anchor cut replaces the
asymptotic end of a broken line with initial monomial $z^m$.  The ordered-chain formula of
Theorem~\ref{thm:seedinversion} expresses each initial coefficient as an
alternating sum of anchored values weighted by bend chains.

\subsubsection*{Further results}
The \emph{good-basis period functional} of a diagonal is the linear functional
given by the coefficient, on that diagonal, of the oscillatory periods of
\cite[Thm.~0.3]{GKT} (Lemma~\ref{lem:periods}).  We call the
$(J,\mathbf d)$-diagonal \emph{narrow} if $r(J)\le r-2$ and $s(J)\le s-2$, and
\emph{broad} if $r(J)=r-1$ or $s(J)=s-1$: as for marked points in
\cite[\S1.1]{GKT}, its residue twist $(r(J),s(J))$ is then Neveu--Schwarz for
both bundles, respectively Ramond for one of them.  The good-basis periods
vanish on the broad diagonals, which are therefore invisible to them.

\begin{theorem}[Effective determination; \S\ref{sec:diagonal}]
\label{thm:effective}
(i) On each narrow $(J,\mathbf d)$-diagonal the wall vectors span exactly
the kernel of the good-basis period functional.  (ii) On each
$(J,\mathbf d)$-diagonal the chamber-index conditions of \cite{GKT} amount to
one linear condition, with the Gamma-function weights of \eqref{eq:Apoly} and
with right-hand side determined by strictly lower order.  (iii) Together with
the characterization of $\Om(I)$ as the set of solutions of the chamber-index
conditions \cite[Cor.~5.14]{GKT}, this determines all open FJRW invariants of
$x^{r}+y^{s}$, up to the torsor action, from the normalizations alone.  (iv) On
every diagonal that carries a balanced coefficient the prescribed value of
$\Ascr$ is $0$ or, for a single marking with descendent exponent $d_1$,
$(-1)^{d_1}$; the closed extended invariants occur only as the values of
$\Ascr$ on diagonals without balanced coefficients \cite[Cor.~5.11(3)]{GKT}.
\end{theorem}

Theorem~\ref{thm:effective} says that the boundary-ray action, together with
the period functional, determines the open theory up to that action.  That
the chamber-index conditions determine the invariants up to the action
follows from the transitivity of \cite[Thm.~4.27]{GKT} together with
\cite[Cor.~5.14]{GKT}, once the Lie algebra is enlarged
(Lemma~\ref{lem:enlarged}).  What is added is the explicit identification of
the wall vectors on each diagonal. On a narrow diagonal they span exactly
the kernel of the good-basis period functional, which also follows from
Steps~1 and~3 of \cite[proof of Thm.~4.27]{GKT}. On a broad diagonal, where
the good-basis periods vanish and \cite[Thm.~4.9]{GKT} does not apply, the
Laplace functional of Lemma~\ref{lem:allresidue} replaces the period formula.
Lemma~\ref{lem:allresidue} shows, by integration by parts, that every
wall term preserves the Laplace integral of $e^{W/\hbar}$ over the positive
quadrant, whose coefficients are fixed nonzero multiples of the polynomials
$\Ascr$ for all residues, the broad ones included.  Besides this lemma, the
proof of Theorem~\ref{thm:effective} has two ingredients.  Each wall vector is
annihilated by the period functional, because the period moments obey the
recurrence of Lemma~\ref{lem:periods}, which is the form that
$\Gamma(z+1)=z\Gamma(z)$ takes on the good-basis cycles.  The $N$ wall vectors
of a diagonal with $N+1$ coefficients are independent and lie in the kernel of
the linear part of $\Ascr$, a nonzero functional on these coefficients, so
they span that kernel, which has dimension $N$.

The last theorem, which is also in \cite{note}, is expressed through the central charge
$\chat_{W}=2-\tfrac2r-\tfrac2s$.  We call the pair $(r,s)$ \emph{hyperbolic} if
$\chat_W>1$, \emph{parabolic} if $\chat_W=1$, and \emph{simple} if $\chat_W<1$.
As in \cite{note}, \emph{hyperbolic} refers to the sign of
$rs-2r-2s=rs(\chat_W-1)$, the sign that makes the triangle group $(2,r,s)$
hyperbolic, and not to Arnold's hyperbolic singularities $T_{p,q,r}$.  The
\emph{threshold} is the value $\chat_W=1$.  The \emph{wall terms over a
monomial} $u_{J,\mathbf d}$ are the wall terms $u_{J,\mathbf d}X_k$ of
$\widetilde\gfrak^{\,r,s}$ with coefficient $u_{J,\mathbf d}$, including the
torus term $u_{J,\mathbf d}X_{0,0}$, $X_{0,0}=x\partial_x-y\partial_y$, when it
is present; each is the wall term of a critical graph over $(J,\mathbf d)$
(Proposition~\ref{prop:shift}).

\begin{theorem}[Weight trichotomy and the threshold; \S\ref{sec:algebra}]
\label{thm:trichotomy}
In the enlarged algebra $\widetilde\gfrak^{\,r,s}$ defined above, the
following hold.
\begin{enumerate}
\item[(i)] For a nonempty marking multiset $J$, wall terms over the
monomial $u_{J,\mathbf d}$ exist if and only if $N(J,\mathbf d)\ge1$, which
forces $w(J,\mathbf d)\ge0$.
\item[(ii)] The wall terms of nonzero bidegree are exactly those of positive
weight; the others are the weight-zero (torus) terms. At $w=0$ the
constraint \eqref{eq:dimconstraint} admits only the marginal torus bidegree
$(0,0)$, and at $w<0$ it admits nothing. Consequently a monomial of
nonpositive weight carries no wall term of nonzero bidegree of its own.
\item[(iii)] The central charge $\chat_{W}$ governs a trichotomy:
positive-weight $d=0$ wall-crossing exists if and only if $\chat_{W}>1$; at
$\chat_{W}=1$ primary wall-crossing exists but is carried entirely by the
marginal torus direction; for $\chat_{W}<1$ there is no primary wall-crossing
at all.
\end{enumerate}
\end{theorem}

Positive weight of the monomial is necessary but not sufficient for a
positive-weight wall term: at $r=s=5$ the monomial $t_{3,3}$ has
$w=5$ but $N=0$, hence no wall term.  The qualifier \emph{of its own} is needed,
since higher-order transport from divisor monomials may reach the coefficients
of a monomial of nonpositive weight.  The threshold is a weight count: a
primary insertion contributes at most $rs(\chat_W-1)$ to the Euler weight, with
equality at the maximal twist $(r-2,s-2)$.  Conversely, when $\chat_W>1$, the
monomial of sufficiently many insertions of maximal twist has $N\ge1$ and
positive weight, and so carries positive-weight wall terms.

\subsection{Comparison with the literature}
Our computations bear on three points in \cite{GKT,GKTsurvey}, recorded here
and, with the supporting derivations, in Remark~\ref{rem:torus}, Proposition~\ref{prop:nontransitive}, Remark~\ref{rem:ex55} and Lemmas~\ref{lem:allresidue} and~\ref{lem:enlarged}.
Theorem, equation and section numbers for \cite{GKT} refer to
arXiv:2203.02435v3, which is the text quoted throughout.

First, as the companion note shows \cite[Rem.~3.1 and Ex.~3.2]{note}, the
marginal torus summands adjoined in \S\ref{ssec:results} are
needed for the transitivity of \cite[Thm.~4.27]{GKT}: with the Lie algebra of
\cite[Def.~4.23]{GKT}, the action is not transitive, for instance, for $r=s=4$
and two labels of twist $(2,2)$ (Proposition~\ref{prop:nontransitive}), and
Remark~\ref{rem:torus} locates where the definitions and proofs of \cite{GKT}
treat this direction.
On the primary part of the parabolic cases the enlarged group is the elliptic
wall-crossing group of Maher's thesis \cite[\S5.2.2]{Maher}: the primary
wall-crossing group that Maher determines from open Saito potentials for
simple elliptic singularities, generated by the products of
$x\partial_x-y\partial_y$ with those positive powers of the marginal
deformation parameter allowed by \cite[(5.2.8)]{Maher} (see
\cite[Cor.~5.2.7, Prop.~5.2.10]{Maher}).

Second, there is an apparent discrepancy in a worked example of the survey.
Imposing the open mirror theorem \cite[Thm.~0.3]{GKT} directly on the
oscillatory expansion of the potential, which is the mechanism of
\cite[Thm.~0.2]{GKT}, gives the two-insertion values \eqref{eq:seed}; at $r=s=4$ this
yields $-\tfrac14$ for the sum in \cite[Ex.~5.5]{GKTsurvey}, where the value
$-\tfrac18$ is stated.  Remark~\ref{rem:ex55} traces this apparent
factor-of-two discrepancy to the count of decompositions into two blocks in
the quadratic part of $\Ascr$, in the numerical evaluation of
\cite[Ex.~5.5]{GKTsurvey}, and confirms the count
independently through \cite[Cor.~5.11(3)]{GKT}; the open mirror theorem itself
is not affected.

Third, two points concern the proof of \cite[Thm.~4.27]{GKT}.  Its
preservation step \cite[proof of Thm.~4.27, Step~1]{GKT} uses the period
formula of \cite[Thm.~4.9]{GKT}, which is stated for residues at most $r-2$
and $s-2$; Lemma~\ref{lem:allresidue} establishes the preservation also for
the broad residues $r(J)=r-1$ or $s(J)=s-1$, as does the proof of
\cite[Prop.~1.4]{note}.  And the argument of
\cite[Thm.~4.27, Step~2]{GKT} establishes freeness of the action, which is
stated there as faithfulness; freeness is the form used here
(Lemma~\ref{lem:enlarged}).

\subsection{Method}
The argument combines the grading and torsor action of the GKT wall-crossing
group with ordered factorization in a strictly convex cone, the
Kontsevich--Soibelman ordered-product lemma over the boundary rays (see the
discussion after Theorem~\ref{thm:presentation}).  The period
relations determine the transverse coordinate on each diagonal, with the
Laplace functional of Lemma~\ref{lem:allresidue} on the broad diagonals.  The
GKT extension lemma realizes each prescribed finite
sequence of group factors by canonical homotopies.
Theorem~\ref{thm:trichotomy} is proved here (\S\ref{sec:algebra}) and also in
\cite{note}, which proves in addition the central-charge threshold for the
dependence of the primary invariants on boundary conditions, and the
first-wall count, the least number of primary interior insertions that admits
a critical graph.  The exact calculations for $x^5+y^5$ in
Section~\ref{sec:cert} and Appendix~\ref{sec:x5y5} are consistency checks and
enter no proof; the programs are available at
\url{https://github.com/bwuebben/open-fjrw-scattering}.

\section{The wall-crossing algebra of critical graphs}
\label{sec:algebra}

We recall the objects of \cite[\S2.4, \S4.3]{GKT} in the notation used
throughout, with the internal and boundary markings of \S\ref{ssec:theory}.
For a multiset $J$ of internal
twists with descendents $\mathbf d$, set $r(J)\equiv\sum_{j}a_{j}$
$(\mathrm{mod}\ r)$, $s(J)\equiv\sum_{j}b_{j}$ $(\mathrm{mod}\ s)$, and
$m(J,\mathbf d)$ as in \eqref{eq:dimconstraint}.

\begin{definition}[{\cite[Not.~2.40]{GKT}}]
\label{def:diagonal}
Write $\sum_{j}a_{j}=r(J)+\ell_{1}r$, $\sum_{j}b_{j}=s(J)+\ell_{2}s$ and
$N=N(J,\mathbf d)=\ell_{1}+\ell_{2}-|J|+1+\sum_{j}d_{j}$. The \emph{balanced}
graphs over $(J,\mathbf d)$ are $\Gamma_{J,p}$, $0\le p\le N$, with boundary
degrees $\beta=(r(J)+pr,\ s(J)+(N{-}p)s)$ and one $\sigma_{12}$-point; the
\emph{critical} graphs are $\Lambda_{J,p}$, $1\le p\le N$, with boundary degrees
$(r(J)+(p{-}1)r+1,\ s(J)+(N{-}p)s+1)$ and no $\sigma_{12}$-point. We refer to the
set of balanced coefficients over a fixed $(J,\mathbf d)$ as its \emph{diagonal}.
The balanced range is empty for $N<0$, and the critical range is empty for
$N<1$. By \cite[Prop.~2.31]{GKT}, the diagonal is the set of boundary degrees
$\beta\in\ZZ_{\ge0}^2$ with $s\beta_{1}+r\beta_{2}=m(J,\mathbf d)$,
$\beta_1\equiv r(J)\ (\mathrm{mod}\ r)$ and
$\beta_2\equiv s(J)\ (\mathrm{mod}\ s)$.  A diagonal thus lies on a segment
parallel to $(r,-s)$ in the boundary-degree plane; it is not the $(1,1)$-ray of
that plane.
\end{definition}

A system of open invariants bounded by $I$ is a vector
$\nu=(\nu_{\Gamma,\mathbf d})$ indexed by the balanced graphs $\Gamma$ with
internal labels in $I$ and their descendent vectors $\mathbf d$
\cite[(0.3)]{GKT}.  For $J\subseteq I$ and $\mathbf d$, the chamber-index
polynomial of \cite[Notation~3.42]{GKT} is
\begin{equation}
\label{eq:Apoly}
\Ascr(J,\mathbf d,\nu)=\sum_{h=1}^{|J|}\frac1{h!}
\sum_{J_1\sqcup\cdots\sqcup J_h=J}\ \sum_{\Theta_1,\dots,\Theta_h}
\frac{\Gamma\bigl(\frac{1+\sum_i \beta_1(\Theta_i)}{r}\bigr)
      \Gamma\bigl(\frac{1+\sum_i \beta_2(\Theta_i)}{s}\bigr)}
     {\Gamma\bigl(\frac{1+r(J)}{r}\bigr)\Gamma\bigl(\frac{1+s(J)}{s}\bigr)}
\prod_{i=1}^h\nu_{\Theta_i,\mathbf d|_{J_i}},
\end{equation}
the sum over ordered decompositions of $J$ into nonempty blocks and over
balanced graphs $\Theta_i$ (so written to keep $\Gamma$ for the Gamma
function) with internal labels $J_i$, where
$\beta(\Theta_i)=(\beta_1(\Theta_i),\beta_2(\Theta_i))$ is the boundary degree of
$\Theta_i$, its numbers of boundary markings of types $\sigma_1$ and $\sigma_2$.  A
\emph{chamber index} is a system with $\nu_{\Gamma,\mathbf 0}=1$ when $\Gamma$
has one internal marking, $\nu_{\Gamma,\mathbf 0}=-1$ when it has none,
$\Ascr(J,\mathbf d,\nu)=(-1)^d$ when $J=\{j\}$ and $\mathbf d=(d)$, and
$\Ascr(J,\mathbf d,\nu)=0$ when $|J|\ge2$ and
$m(J,\mathbf d)\ge s\,r(J)+r\,s(J)$ \cite[Definition~3.47]{GKT}; we write
$\Upsilon_{r,s}(I)$ for the set of chamber indices bounded by $I$.  By
\cite[Corollary~5.14]{GKT}, $\Om(I)=\Upsilon_{r,s}(I)$: the chamber indices are
exactly the systems realized by canonical boundary conditions, the
\emph{realizable} systems of \S\ref{ssec:theory}.
Since $m(J,\mathbf d)=s\,r(J)+r\,s(J)+rs\,N(J,\mathbf d)$, the inequality in
the last condition says $N(J,\mathbf d)\ge0$.

\begin{proposition}[The primitive-form shift]
\label{prop:shift}
The Lie bidegree of the critical graph $\Lambda_{J,p}$, in the sense of
\eqref{eq:dimconstraint}, is its boundary degree minus $(1,1)$
(\cite[Prop.~2.39]{GKT}); equivalently, the vector field of the wall term attached
to $\Lambda$ is the Hamiltonian vector field, for the primitive form $dx\wedge dy$, of the boundary
monomial $x^{\#\sigma_{1}}y^{\#\sigma_{2}}$ of the critical graph. With
$g_{\alpha}=z^{\alpha}(\alpha_{2}x\partial_{x}-\alpha_{1}y\partial_{y})$ the corresponding vector
field of the tropical vertex (Hamiltonian for $d\log x\wedge d\log y$), one has
$X_{k_{1},k_{2}}=(xy)^{-1}g_{(k_{1}+1,k_{2}+1)}$, and the
wall-crossings are the rational-power maps of
\cite[Lem.~4.21]{GKT}. The shift is the removal of the $\sigma_{12}$-insertion.
The algebra is $\ZZ^2$-graded by $(k_1,k_2)$; the ray assignment uses
$(k_1,k_2)+(1,1)$, which is not additive under the bracket.
\end{proposition}

\begin{proof}
The first statement is \cite[Prop.~2.39(b)]{GKT} read against
Definition~\ref{def:diagonal}. Since $dx\wedge dy=xy\,d\log x\wedge d\log y$,
the Hamiltonian vector field of a function for $dx\wedge dy$ is $(xy)^{-1}$
times its Hamiltonian vector field for $d\log x\wedge d\log y$; for
$x^{k_1+1}y^{k_2+1}$ this is the dictionary $X_{k}=(xy)^{-1}g_{k+(1,1)}$, and
the explicit wall maps are \cite[Lem.~4.21]{GKT}.  The grading follows from
$[X_{k},X_{k'}]\in\QQ\,X_{k+k'}$ (\S\ref{sec:presentation}).
\end{proof}

\begin{proof}[Proof of Theorem~\ref{thm:trichotomy}]
The constraint \eqref{eq:dimconstraint} reads $sk_{1}+rk_{2}=w(J,\mathbf d)$ with
$k_{1},k_{2}\ge0$: for $w<0$ there are no solutions, for $w=0$ only $(0,0)$, and
for $w>0$ the $N$ solutions of Definition~\ref{def:diagonal}. A twist $(a,b)$
contributes $sa+rb-rs\le s(r{-}2)+r(s{-}2)-rs=rs(\chat_{W}-1)$ to the weight, with
equality at the maximal twist; descendents contribute $rs\,d_{j}>0$. Hence
positive-weight $d=0$ monomials exist if and only if $\chat_{W}>1$, and for
$\chat_{W}=1$ every $d=0$ monomial has $w\le0$ with equality exactly on products of
marginal twists; whenever such an equality monomial is wall-bearing, its unique
direction is the torus one. For the
existence of wall terms (not merely of positive-weight monomials) when
$\chat_{W}>1$: take $n$ interior markings of maximal twist $(r{-}2,s{-}2)$;
then $N=n-\lceil2n/r\rceil-\lceil2n/s\rceil+1\ge1$ once
$n(1-\tfrac2r-\tfrac2s)\ge2$, while the weight is $n\,rs(\chat_{W}-1)>0$, so
the resulting critical graphs carry wall terms; this is the maximal-twist
construction of \cite{note}. At $\chat_{W}=1$, those analogous multisets for
which $N\ge1$ are critical and marginal, realizing the torus wall-crossing; for
$\chat_{W}<1$ and any primary $J$ with $|J|=n$,
$\ell_{1}+\ell_{2}\le n\chat_{W}<n$ forces $N\le0$: no primary critical graphs
at all. The qualifier in the theorem is necessary: acting on a lower divisor
coefficient multiplies deformation monomials, so nonlinear transport can
contribute to a nonpositive-weight target even though no wall term is supported there.
\end{proof}

\begin{remark}[The marginal torus direction]
\label{rem:torus}
At a wall-bearing marginal monomial ($w=0$ and $N\ge1$; necessarily
$r(J)=s(J)=0$ and $N=1$) the critical graph exists (for $x^{4}+y^{4}$ it
is precisely the disk of \cite[Fig.~9, Ex.~4.9]{GKTsurvey}), but its Lie
bidegree is $(0,0)$, the direction $x\partial_{x}-y\partial_{y}$, which is not
included in the explicit description \cite[(4.16), (4.19)]{GKT}, although the
intrinsic description of the group in \cite[\S0.4 and Rem.~4.24]{GKT} admits
it.  With coefficients in the ideal generated by the deformation variables,
the exponential of this direction converges adically.  The proofs of
\cite[Thm.~5.9]{GKT} and of \cite[Thm.~5.13]{GKT} (which gives
\cite[Cor.~5.14]{GKT}) realize the torus motion.  In
the proof of \cite[Thm.~5.9]{GKT}, the jump sets $Q_{0}(A,\mathbf d)$ of the
jump formula (\S\ref{sec:presentation}) impose
only the congruences and \eqref{eq:dimconstraint}, which admit the bidegree
$(0,0)$ when $r(A)=s(A)=0$.  In the
proof of \cite[Thm.~5.13]{GKT} (which gives \cite[Cor.~5.14]{GKT}), the
realizability step modifies the
multisection of a critical graph $\Lambda$ by a homotopy
\cite[Lem.~5.17]{GKT}.  This moves the two adjacent balanced coefficients by
$h_\Lambda\cdot((-1)^{\beta_{2}(\Lambda)}s\,\beta_{1}(\Lambda),
(-1)^{\beta_{2}(\Lambda)+1}r\,\beta_{2}(\Lambda))$, where $\beta(\Lambda)$ is the
boundary degree of $\Lambda$ and $h_\Lambda\in\QQ$ is the signed number of
zeros of that homotopy, the number $h_{p,\mathbf d}$ of
\cite[(5.52)--(5.54)]{GKT}.  On a marginal diagonal this vector is exactly the
antisymmetric torus motion.
We therefore work with the enlarged algebra
\[
\widetilde\gfrak^{\,r,s}=\gfrak^{r,s}\ \oplus\
\bigoplus_{\substack{\varnothing\ne J,\ w(J,\mathbf d)=0,\ N(J,\mathbf d)\ge1\\
r(J)=s(J)=0}}
\QQ\,u_{J,\mathbf d}\,(x\partial_{x}-y\partial_{y}),
\]
whose new summands commute with one another and preserve the bidegree of
every wall term. This is a Lie algebra: with $X_{0,0}=x\partial_x-y\partial_y$,
$[uX_{0,0},vX_{k}]=(k_1-k_2)uvX_{k}$ and $[uX_{0,0},vX_{0,0}]=0$; since a
marginal coefficient
$u$ has weight and residues zero, $uv$ satisfies the same wall-term
constraints as $v$. In the primary algebra nothing is added when $\chat_{W}<1$; for the parabolic pairs the primary
group is purely torus, matching the elliptic wall-crossing group of Maher
\cite[\S5.2.2]{Maher}; for $\chat_{W}>1$ the torus summands sit over the marginal
submonomials.  Without the torus summands transitivity fails: for $r=s=4$
and two labels of twist $(2,2)$, the chamber-index conditions leave a
one-parameter set of values for the coefficients of the product of the two
primary label variables, the set exhibited in \cite[Ex.~4.9]{GKTsurvey}, and no
generator of $\gfrak^{4,4}$ changes them
(Proposition~\ref{prop:nontransitive}). All statements below are over $\widetilde\gfrak^{\,r,s}$ and its
group $\widetilde G^{r,s}=\exp(\widetilde\gfrak^{\,r,s})$; the next two lemmas
supply the torsor properties over the enlargement, which over the symmetric
ring are \cite[Prop.~1.4]{note}.
\end{remark}

\begin{lemma}[All-residue $\Gamma$-functionals]
\label{lem:allresidue}
For every nonempty labeled marking set $J$ and descendent vector $\mathbf d$,
the polynomial $\Ascr(J,\mathbf d,\nu)$ is preserved by every wall term of
$\widetilde\gfrak^{\,r,s}$, including when $r(J)=r-1$ or $s(J)=s-1$.  (The
preservation step of \cite[proof of Thm.~4.27, Step~1]{GKT} uses the period
formula of \cite[Thm.~4.9]{GKT}, which is stated for residues at most $r-2$
and $s-2$.)
\end{lemma}

\begin{proof}
Take $\hbar<0$ and expand, coefficientwise in the deformation variables, the
positive-quadrant Laplace integral
\[
 \mathcal L_\hbar(W)=
 \int_{0}^{\infty}\!\int_{0}^{\infty} e^{W/\hbar}\,dx\,dy .
\]
The monomial moments are
\[
 \int_{0}^{\infty}x^q e^{x^r/\hbar}\,dx
 =\frac1r(-\hbar)^{(q+1)/r}
   \Gamma\!\left(\frac{q+1}{r}\right),
\]
and similarly for $y$.  Therefore the coefficient of
$u_{J,\mathbf d}$ in $\mathcal L_\hbar(W^\nu)$ is a fixed nonzero power of
$-\hbar$, multiplied by fixed nonzero Gamma factors depending only on
$r(J),s(J)$, times $\Ascr(J,\mathbf d,\nu)$.  This is the termwise
calculation in the proof of \cite[Thm.~4.9]{GKT}, before the good-basis
Kronecker factors are imposed.  Unlike those good-basis periods, the displayed
Laplace functional is nonzero also for residues $r-1$ and $s-1$.

Let $\eta=uX_{k}$ be any wall term of the enlarged algebra.  The vector field $X_{k}$ is
divergence-free and tangent to both coordinate axes.  Hence, coefficientwise,
\[
 \frac{d}{dT}\bigg|_{T=0}\mathcal L_\hbar(e^{T\eta}W)
 =\int_{[0,\infty)^2}\eta(e^{W/\hbar})\,dx\,dy=0:
\]
integration by parts has no contribution on the axes, and the exponential
decay removes the boundary at infinity.  The same identity after replacing
$W$ by $e^{T\eta}W$ shows that the full exponential action preserves
$\mathcal L_\hbar$, and comparison of deformation coefficients proves the
claim.  Equivalently, the calculation is the Gamma recurrence
$\Gamma(z+1)=z\Gamma(z)$ on every residue class.
\end{proof}

\begin{lemma}[The enlarged torsor]
\label{lem:enlarged}
The action of $\widetilde G^{r,s}$ on potentials over the labeled ring
$A_I$ of \cite[Def.~4.6]{GKT} preserves $\Om(I)$ and is \emph{free} and
transitive on it, so that $\Om(I)$ is a torsor in the strict sense, and the
jump automorphisms of \cite[Thm.~5.9]{GKT} lie in $\widetilde G^{r,s}$.  The
same holds for the symmetric systems over the symmetric ring
$A_{I,\mathrm{sym}}$.
\end{lemma}

\begin{proof}
Membership of the jumps: the jump sets $Q_{0}(A,\mathbf d)$ of the jump
formula (\S\ref{sec:presentation}) impose only the congruences and
\eqref{eq:dimconstraint}, which admit the bidegree $(0,0)$ when
$r(A)=s(A)=0$, so the jump elements are sums of positive-weight and
weight-zero (torus) wall terms.

Preservation: the coefficient of every wall term has positive deformation
degree, so its exponential converges in that filtration.  Lemma~\ref{lem:allresidue}
preserves all the $\Ascr$-conditions of \cite[Def.~3.47]{GKT}.  It also
preserves the normalization there: a primary critical monomial has at least
two markings, while a descendent singleton does not alter the
$\mathbf d=\mathbf0$ normalization.  Thus the enlarged group preserves
$\Upsilon_{r,s}(I)$.

Transitivity: first work modulo a finite divisor-closed set of deformation
monomials and order those monomials by any linear extension of proper
divisibility.  Suppose two realizable systems have already been made equal on
every proper divisor of $u=u_{J,\mathbf d}$.  On the diagonal of $u$, all
$h\ge2$ terms of $\Ascr(J,\mathbf d,-)$ involve proper divisors and hence
already agree.  The difference of the two coefficient vectors therefore lies
in the kernel of the nonzero $h=1$ Gamma functional.  If the diagonal has
$N+1$ balanced coefficients, that kernel has dimension $N$.  In the basis
$\varepsilon_0,\dots,\varepsilon_N$ of the coefficients of the diagonal, the $N$
adjacent wall vectors (computed in \S\ref{sec:diagonal})
\[
 v_p=\lambda_p\varepsilon_p-\mu_p\varepsilon_{p-1},\qquad \lambda_p,\mu_p>0,\qquad 1\le p\le N,
\]
are independent and lie in the kernel by Lemma~\ref{lem:allresidue}; hence
they span it.  On a wall-bearing marginal diagonal $N=1$, and the sole vector is the
adjoined torus vector
$r\varepsilon_1-s\varepsilon_0=X_{0,0}(W_0)$.  A unique combination of
wall terms with coefficient $u$ therefore removes the difference.  Its action
on every nonconstant term of the potential carries a strict multiple of $u$, so
it changes no coefficient on a proper divisor of $u$ and acts on the diagonal
of $u$ only through $W_0$.  Induction proves
transitivity in every finite quotient.  By the freeness proved below, whose
argument applies verbatim modulo the monomials outside a finite
divisor-closed set, the element constructed in each quotient is unique, so
the elements for nested quotients are compatible; since only finitely many
deformation monomials occur below each bound on the degree of the wall terms,
they define an element of the completed group.

The equality $\Om(I)=\Upsilon_{r,s}(I)$ remains the realization theorem
\cite[Thm.~5.13 and Cor.~5.14]{GKT}.  Its inclusion
$\Upsilon_{r,s}(I)\subseteq\Om(I)$ does not require the omitted marginal
transitivity: at equation (5.50) of \cite{GKT}, the two required
$\Ascr$-values agree directly because both are the prescribed value of
\cite[Def.~3.47]{GKT}.  The enlarged action therefore preserves and is
transitive on the realizable systems.

Freeness: let $g=\exp(\eta)$ with $\eta\ne0$, and let $u$ be minimal in the
divisibility order among the monomials in the support of $\eta$ (labeled
monomials over $A_I$, symmetric monomials over $A_{I,\mathrm{sym}}$; the
argument is the same). The change of
the $u$-diagonal of any potential is exactly $\eta_{u}(W_{0})$, $W_{0}=x^{r}+y^{s}$: contributions of
other wall terms, or of higher exponential terms, carry monomials properly
containing a support monomial, which cannot equal $u$ by minimality. And
$\eta_{u}(W_{0})\ne0$. Indeed, on the diagonal belonging to $u$, the $N$
positive-weight wall terms act by the adjacent vectors
$\lambda_{p}\varepsilon_{p}-\mu_{p}\varepsilon_{p-1}$ with $\lambda_{p},\mu_{p}>0$. These
vectors are linearly independent: reading a relation successively in the
$\varepsilon_{N},\varepsilon_{N-1},\ldots,\varepsilon_{1}$ coordinates kills all its
coefficients. On a marginal diagonal there is instead the single nonzero torus
vector $r\varepsilon_{1}-s\varepsilon_{0}$. Thus no nonzero linear combination of the wall terms in
$\eta_u$ can cancel. (Individually, a positive-weight wall term contributes
$r(k_{2}{+}1)x^{k_{1}+r}y^{k_{2}}-s(k_{1}{+}1)x^{k_{1}}y^{k_{2}+s}$, while a
torus term contributes $r\,x^{r}-s\,y^{s}$.) So $g$ moves every potential, which is freeness and not
merely faithfulness. The distinction matters downstream: uniqueness of the
transition element, used in Theorem~\ref{thm:presentation}(i), needs trivial
stabilizers, and faithfulness alone does not supply them for a nonabelian group.
The argument of \cite[Thm.~4.27, Step~2]{GKT} likewise establishes freeness,
though it is stated there as faithfulness. All uses below are through this lemma.
\end{proof}

\begin{proposition}[Transitivity requires the torus summands]
\label{prop:nontransitive}
Let $r=s=4$, and let $I=\{i,j\}$ consist of two labels of twist $(2,2)$, with
primary label variables $u_i,u_j$.  The chamber-index conditions leave a
one-parameter set of values for the coefficients of $u_iu_j$ in the
potential, the set exhibited in \cite[Ex.~4.9]{GKTsurvey}, and no generator of the
Lie algebra $\gfrak^{4,4}$ of \cite[Def.~4.23]{GKT} changes these
coefficients.  Consequently $\exp(\gfrak^{4,4})$ does not act transitively on
$\Om(I)$.
\end{proposition}

\begin{proof}
For $J=\{i,j\}$ and $\mathbf d=\mathbf 0$ one has $r(J)=s(J)=0$,
$\ell_1=\ell_2=1$, $N=1$ and $w=0$, with balanced boundary degrees $(0,4)$ and
$(4,0)$.  The chamber-index condition on this diagonal is the single condition
$\Ascr(J,\mathbf 0,\nu)=0$, whose linear part has two nonzero weights, so it
leaves a one-parameter set of values for the two coefficients.  Since
$u_iu_j$ divides no other monomial of the labeled ring, no other condition
involves these coefficients, so every value in this set occurs in a chamber index, hence
in $\Om(I)$ by \cite[Cor.~5.14]{GKT}.  A generator multiplies the deformation
monomials of the terms on which it acts by its own coefficient, so it can
change these coefficients only if its coefficient is $u_i$, $u_j$ or
$u_iu_j$.  The singletons have $N=0$ and carry no generator, and $u_iu_j$ has
weight zero, so by \eqref{eq:dimconstraint} its only admissible bidegree is
$(0,0)$, which $\gfrak^{4,4}$ does not contain.  Hence every element of
$\exp(\gfrak^{4,4})$ fixes these coefficients, and two realizable systems with
different values in this set lie in different orbits.
\end{proof}

\section{Slope-ordered boundary-ray factorization}
\label{sec:presentation}

Throughout, a \emph{normal-form homotopy} is a homotopy of boundary
conditions supplied by \cite[Lem.~3.40]{GKT}, recalled as
Lemma~\ref{lem:gktextension} at the end of this section: for each subset $A$ of the
internal labels there is one time $T_{A}\in(0,1)$, all times distinct, and
every jump of a critical graph with internal marking set $A$ occurs at
$T_{A}$. The times are auxiliary choices (Remark~\ref{rem:orderinggap}).
We restate the jump formula of \cite[proof of Thm.~5.9]{GKT}.  The potential
$W^{\mathfrak s_T}$ of such a homotopy is constant in the homotopy time $T$ except at the times $T_A$,
and at $T_A$ it changes by
\[
W\longmapsto\exp\Bigl(\sum_{\mathbf d}\sum_{\Lambda}
(-1)^{|A|+k_2(\Lambda)}\,\Delta(\Lambda)\,u_{A,\mathbf d}X_{\Lambda}\Bigr)(W)
\qquad\text{\cite[proof of Thm.~5.9]{GKT}},
\]
the sums over the descendent vectors $\mathbf d$ on $A$ and the critical graphs
$\Lambda$ with internal marking set $A$ and descendent vector $\mathbf d$.
Here $X_\Lambda=X_{k(\Lambda)}$, where $k(\Lambda)=\beta(\Lambda)-(1,1)$ is the
Lie bidegree of $\Lambda$ (Proposition~\ref{prop:shift}) and $k_2(\Lambda)$ its
second component; $u_{A,\mathbf d}=\prod_{j\in A}u_{j,d_j}$ is the monomial of
the internal markings in the labeled ring of \cite[Def.~4.6]{GKT}, so that
$u_{A,\mathbf d}X_\Lambda$ is the wall term of $\Lambda$; and $\Delta(\Lambda)$ is
the signed count of zeros of the homotopy of
$\Lambda$ on $[0,1]\times\mathcal{PM}_\Lambda$ \cite[(5.33)]{GKT}, with
$\mathcal{PM}_\Lambda$ GKT's partial compactification of the moduli space of
$\Lambda$.  These
automorphisms are the \emph{jumps} of the homotopy.  (In \cite{GKT} the sum
over $\Lambda$ runs over the set $Q_0(A,\mathbf d)$ of \cite[Not.~5.7]{GKT},
whose conditions are the congruences and \eqref{eq:dimconstraint} for
$k(\Lambda)$.)

\begin{proof}[Proof of Theorem~\ref{thm:presentation}(i)]
Let $(\mathfrak s_{T})_{T\in[0,1]}$ be a homotopy of boundary conditions
supplied by \cite[Lem.~3.40]{GKT}, with $\mathfrak s_{1}=\mathfrak s_{0}$, jump times
$T_{1}<\dots<T_{M}$, the times of the marking subsets $A_1,\dots,A_M$, and
jump elements $g_{i}\in\widetilde G^{r,s}$ (Lemma~\ref{lem:enlarged}) given by
the jump formula above
(\cite[proof of Thm.~5.9, eqns.~(5.33)--(5.37)]{GKT}). By that proof the potential
$W^{\mathfrak s_{T}}$ is constant between jumps and transforms by $g_{i}$ at $T_{i}$;
around the loop, $W^{\mathfrak s_{0}}=(g_{M}\circ\cdots\circ g_{1})(W^{\mathfrak s_{0}})$, and
freeness (Lemma~\ref{lem:enlarged}) gives $g_{M}\circ\cdots\circ
g_{1}=\mathrm{id}$. For two such paths with fixed endpoints, their total jumps
$g,g'$ satisfy $g'g^{-1}(W^{\mathfrak s_{1}})=W^{\mathfrak s_{1}}$, hence $g'=g$.
\end{proof}

\begin{proof}[Proof of Theorem~\ref{thm:presentation}(ii)]
Let $\mathfrak m$ be the deformation ideal and first work in the primary
quotient
\[
 F^{p}\widetilde\gfrak^{\,r,s}_{0}
 :=\operatorname{span}\{u_{J,\mathbf0}X:\lvert J\rvert\ge p\},
 \qquad
 L_{<K}:=\widetilde\gfrak^{\,r,s}_{0}/F^{K}.
\]
It is a finite-dimensional nilpotent Lie algebra and
$[F^{p},F^{q}]\subseteq F^{p+q}$ (here and below $F^\bullet$ also denotes the
induced filtration on the quotient). For a rational boundary ray $\rho$ let
$L_{\rho}$ be the span of the wall terms $uX_{k_{1},k_{2}}$ for which
$(k_{1}{+}1,k_{2}{+}1)\in\rho$, placing a marginal wall term
$u(x\partial_x-y\partial_y)$ in $\rho=\RR_{\ge0}(1,1)$. Then
$L_{<K}=\bigoplus_{\rho}L_{\rho}$, with only finitely many summands.

Each $L_{\rho}$ is abelian. Indeed,
\[
[X_{k_{1},k_{2}},X_{k'_{1},k'_{2}}]
=\det\bigl((k'_{1}{+}1,k'_{2}{+}1),(k_{1}{+}1,k_{2}{+}1)\bigr)
X_{k_{1}+k'_{1},\,k_{2}+k'_{2}},
\]
so the bracket vanishes when the boundary-ray directions
$(k_1{+}1,k_2{+}1)$ and $(k'_1{+}1,k'_2{+}1)$ are parallel; on the
$(1,1)$-ray the marginal vector field also commutes with every $X_{q,q}$.
The deformation coefficient of a marginal wall term has positive
$\mathfrak m$-degree, so it belongs to the same nilpotent filtration.

Fix any slope order. The multiplication map
\[
\prod_{\rho}^{\longrightarrow}\exp(L_{\rho})
\longrightarrow \exp(L_{<K})
\]
is bijective. Modulo $F^{2}$ this is the direct-sum decomposition above.
Suppose the factors have been found through $F^{n}$. The logarithm of the
residual in $F^{n}/F^{n+1}$ decomposes uniquely by $\rho$; insert these
components into the corresponding factors. Commuting a degree-$n$ correction
past a previously chosen positive-degree term changes the product only in
$F^{n+1}$. This proves existence by induction, as in the
Kontsevich--Soibelman lemma \cite{KS06}, \cite[Thm.~1.4]{GPS10}. At the first
filtration degree
where two ordered factorizations differ, the same direct-sum decomposition
forces all differences to vanish, proving uniqueness. The construction is
compatible as $K$ varies.

For descendents, perform the same construction after quotienting to any finite
divisor-closed set of deformation monomials. Its complement is upward closed,
so this is a Lie-algebra quotient, and the compatible finite constructions give
the asserted coefficientwise factorization.
\end{proof}

\begin{remark}
\label{rem:timeorder}
The proof of \cite[Thm.~5.9]{GKT} produces the total jump in a homotopy-time-ordered factorization, one factor per subset of internal markings, mixing boundary rays;
the proof of \cite[Lem.~3.40]{GKT} shows the jump times and their order are auxiliary
choices. The boundary-ray factorization is a reorganization, involving no
choices, by the boundary rays, whose directions are the Lie bidegrees
shifted by $(1,1)$ as forced by Proposition~\ref{prop:shift}, and not by the
rays of the Lie bidegrees $k$, which are additive under the bracket; this
shift is what makes the boundary-ray factorization $\Dscr^{r,s}$ of
Theorem~\ref{thm:canonical} geometrically meaningful.
\end{remark}

For later reference we recall the parts of the extension lemma of \cite{GKT}
that are used below.  For a graph $\Gamma$, let $\mathcal{PM}_\Gamma$ be as
above, $E_\Gamma(\mathbf d)$ the descendent Witten bundle and $\#\sigma_{12}(\Gamma)$
the number of $\sigma_{12}$-points (written $k_{12}(\Gamma)$ in \cite{GKT}).  A critical graph is connected and
relevant \cite[Def.~1.57]{GKT}, has $\#\sigma_{12}=0$ and $\operatorname{rank}E_\Gamma(\mathbf d)=\dim\mathcal{PM}_\Gamma+1$,
and a connected balanced graph has $\#\sigma_{12}=1$ and rank equal to dimension
\cite[Defs.~2.32--2.33]{GKT}.

\begin{lemma}[{\cite[Lem.~3.40]{GKT}}]
\label{lem:gktextension}
Let $\mathfrak s_0$ and $\mathfrak s_1$ be families of canonical multisections
bounded by $I$ that are transverse to the zero section.  There are a family
$H^*=(H^\Gamma)_\Gamma$ of simple canonical homotopies bounded by $I$ from
$\mathfrak s_0$ to $\mathfrak s_1$ and distinct times $T_A\in(0,1)$, one for
each $A\subseteq I$, such that the following hold for every smooth graded
$W$-spin graph $\Gamma$ with only open vertices and $I(\Gamma)\subseteq I$
and every descendent vector $\mathbf d$.
\begin{enumerate}
\item[(1)] $H^\Gamma$ vanishes nowhere if $\Gamma$ is irrelevant in the sense
of \cite[Def.~1.57]{GKT}, if $\Gamma$ is connected with
$\operatorname{rank}E_\Gamma(\mathbf d)\ge\dim\mathcal{PM}_\Gamma+2$, or if
two connected components of $\Gamma$ have distinct internal marking sets and
each has rank one more than its dimension and at most two
$\sigma_{12}$-points.
\item[(2)] If $\Gamma$ is connected and either (a)
$\operatorname{rank}E_\Gamma(\mathbf d)=\dim\mathcal{PM}_\Gamma+1$ and
$\#\sigma_{12}(\Gamma)\le2$, or (b)
$\operatorname{rank}E_\Gamma(\mathbf d)=\dim\mathcal{PM}_\Gamma$ and
$\#\sigma_{12}(\Gamma)\le1$, then the restriction of $H^\Gamma$ to every stratum of
$[0,1]\times\mathcal{PM}_\Gamma$ is transverse to the zero section.
\item[(3)] Let $\Gamma$ be connected.  (a) In case (2)(a), $H^\Gamma$ has no
zeros on $([0,1]\setminus\{T_{I(\Gamma)}\})\times\mathcal{PM}_\Gamma$.  (b) In
case (2)(b), for every $A\subseteq I$ the restriction of $H^\Gamma$ to
$\{T_A\}\times(\mathcal{PM}_\Gamma\setminus\partial\mathcal{PM}_\Gamma)$ is
transverse to the zero section.
\end{enumerate}
\end{lemma}

By (3)(a), every zero of the homotopy of a critical graph with internal
marking set $A$ occurs at the time $T_A$, which is the normal form used in
this section; by (3)(b), the slices of the homotopies of balanced graphs at
the times $T_A$ are transverse on
$\mathcal{PM}_\Gamma\setminus\partial\mathcal{PM}_\Gamma$.

\section{Period relations on each diagonal}
\label{sec:diagonal}

Two structural facts from \cite{GKT} organize this section. The open
topological recursion factorizes through a
cohomological-field-theory node splitting: in
\cite[Lem.~6.4]{GKT} the distinguished insertion splits off a closed extended
invariant and continues as an open piece carrying the socle-dual twist
$(a,b)\mapsto(r{-}2{-}a,s{-}2{-}b)$ of the Jacobi ring
$\CC[x,y]/(x^{r-1},y^{s-1})$: the node is the residue pairing. And the
$\Gamma$-weights in the invariant combinations $\Ascr$ are the periods of
$x^{r}+y^{s}$: $\int_{0}^{\infty}e^{-x^{r}}x^{\beta_1}dx=\tfrac1r\Gamma(\tfrac{1+\beta_1}r)$,
so the weight attached to a boundary degree $\beta$ is the period normalization of the
primitive form. The good-basis period moments themselves are explicit:

\begin{lemma}[Periods]
\label{lem:periods}
For the good-basis cycles $\Xi_{a}$ of $x^{r}$
($\int_{\Xi_{a}}x^{a'}e^{x^{r}/\hbar}=\delta_{aa'}$, $0\le a,a'\le r-2$), the
moments $P_{a}(p)=\int_{\Xi_{a}}x^{p}e^{x^{r}/\hbar}$ satisfy
$P_{a}(p)=\delta_{ap}$ for $p\le r{-}2$, $P_{a}(r{-}1)=0$, and
$P_{a}(p)=-\tfrac{\hbar(p-r+1)}{r}P_{a}(p-r)$.
\end{lemma}

\begin{proof}
Total derivatives integrate to zero on $\Xi_{a}$:
$0=\int_{\Xi_{a}}\tfrac{d}{dx}(x^{q}e^{x^{r}/\hbar})
=q\,P_{a}(q{-}1)+\tfrac r\hbar P_{a}(q{+}r{-}1)$.
\end{proof}

\begin{proof}[Proof of Theorem~\ref{thm:effective}]
\emph{The case $N=0$.}
If $N=0$, there is no wall vector and the nonzero $h=1$ weight in
$\Ascr$ determines the sole coefficient modulo lower monomials, so the claim is
immediate.

\emph{Wall-bearing diagonals.}
Fix now a wall-bearing diagonal $(J,\mathbf d)$, of positive or
marginal weight, with coefficients $\nu_{p}$, $0\le p\le N$, as in
Definition~\ref{def:diagonal}. The wall term of
$\Lambda_{J,p}$ acts on $W_{0}=x^{r}+y^{s}$ by
\begin{equation}
\label{eq:wallvector}
u_{J,\mathbf d}X_{k_{1},k_{2}}(W_{0})
=u_{J,\mathbf d}\bigl(r(k_{2}{+}1)\,x^{k_{1}+r}y^{k_{2}}
-s(k_{1}{+}1)\,x^{k_{1}}y^{k_{2}+s}\bigr),
\end{equation}
which lands on the monomials of two adjacent balanced coefficients of the same
diagonal: in the basis
$\varepsilon_{0},\dots,\varepsilon_{N}$ of coefficients, it is the wall vector
$v_{p}=\lambda_{p}\varepsilon_{p}-\mu_{p}\varepsilon_{p-1}$ with
$\lambda_{p}=r(k_{2}{+}1)=r\,\beta_{2}(\Lambda_{J,p})>0$ and
$\mu_{p}=s(k_{1}{+}1)=s\,\beta_{1}(\Lambda_{J,p})>0$, where
$\beta(\Lambda_{J,p})=k+(1,1)$ is the boundary degree of $\Lambda_{J,p}$.
(These are exactly the geometric jump coefficients of \cite[(5.52)--(5.54)]{GKT}.) The $v_{p}$ are triangular, hence
independent, and each is annihilated by the period functional: by
Lemma~\ref{lem:periods}, $P_{a}(k_{1}{+}r)P_{b}(k_{2})
=-\tfrac{\hbar(k_{1}+1)}{r}P_{a}(k_{1})P_{b}(k_{2})$ cancels against the second
term of \eqref{eq:wallvector}.

\emph{Narrow diagonals.}
On a narrow diagonal the period functional is nonzero on each $\varepsilon_{p}$, so its kernel has
dimension $N$ and coincides with $\mathrm{span}(v_{p})$. By
\cite[Thm.~4.9]{GKT} the coefficient of $u_{J,\mathbf d}$ in
$\int_{\Xi_{(r(J),s(J))}}e^{W^{\nu}/\hbar}\,dx\wedge dy$ equals
$(-1)^{|J|}(-\hbar)^{-d(J,\mathbf d)-2}\Ascr(J,\mathbf d,\nu)$, with
$d(J,\mathbf d)=\bigl(s\,r(J)+r\,s(J)-m(J,\mathbf d)\bigr)/rs-1$
\cite[Def.~4.8]{GKT}.  The $h=1$ part of $\Ascr$
($h$ the number of blocks in \eqref{eq:Apoly}) is
the weighted sum $\sum_{p}\omega_{p}\nu_{p}$ over the diagonal, with
$\omega_p$ the $h=1$ weights of \eqref{eq:Apoly}, and its
$h\ge2$ parts involve only strictly smaller monomials. The chamber-index
conditions \cite[Def.~3.47]{GKT} prescribe the value of $\Ascr$.

\emph{Broad diagonals.}
On broad
diagonals ($r(J)=r-1$ or $s(J)=s-1$) the good-basis periods vanish identically
on the diagonal (Lemma~\ref{lem:periods}), but the $h=1$ weights of $\Ascr$
remain nonzero there (their arguments $(1+\beta_1)/r$ and $(1+\beta_2)/s$ are
positive, and $\Gamma$ does not vanish there), and $\mathrm{span}(v_{p})$ lies in the
kernel of the weighted functional modulo strictly smaller monomials, by
Lemma~\ref{lem:allresidue}; so the same
one-linear-condition structure holds, imposed directly by the conditions,
which by \cite[Cor.~5.14]{GKT} characterize $\Om(I)$ completely. Ordering monomials by degree and descendents determines every
invariant, up to the span of the $v_{p}$ (the torsor directions), proving the
theorem.  The prescribed values of $\Ascr$ on diagonals with balanced
coefficients are $0$ or, for a single marking with descendent exponent
$d_1$, $(-1)^{d_1}$, because
$m(J,\mathbf d)=s\,r(J)+r\,s(J)+rs\,N(J,\mathbf d)$ gives
$d(J,\mathbf d)=-N(J,\mathbf d)-1<0$ there \cite[Def.~4.8]{GKT}; this proves
the last statement.  The open topological recursion \cite[Thm.~0.5]{GKT}
then relates these invariants to the closed extended invariants; it is not
among the conditions.
\end{proof}

\begin{remark}[The first values over the twist $(3,3)$]
\label{rem:tower}
For $x^{5}+y^{5}$ the conditions produce, order by order over the
$t_{3,3}$-monomials, the one-insertion normalization
$\langle\tau_{0}^{(3,3)}\sigma_{1}^{3}\sigma_{2}^{3}\sigma_{12}\rangle=1$; the
first-wall relation \eqref{eq:firstwallrel} at $t_{3,3}^{2}$; the absolute value
$\tfrac15$ of \eqref{eq:newinvariants} at $t_{3,3}^{3}$ (a diagonal with $N=0$,
hence no wall terms, and $r(J)=4=r{-}1$, hence broad and invisible to the good-basis
periods, so the condition that determines it is not a period condition); and at $t_{3,3}^{4}$
a relation quadratic in the first-wall invariants, giving
$\nu^{\min}_{(2,7)}=-\tfrac{12}{25}$ in the minimal system of
Theorem~\ref{thm:canonical}. The descendent singleton at $d=1$ gives
$\langle\tau_{1}^{(3,3)}\sigma_{1}^{8}\sigma_{2}^{3}\sigma_{12}\rangle
+\langle\tau_{1}^{(3,3)}\sigma_{1}^{3}\sigma_{2}^{8}\sigma_{12}\rangle
=-\tfrac54$: on its diagonal ($N=1$, balanced boundary degrees $(3,8)$ and $(8,3)$)
the condition $\Ascr=(-1)^1$ has only its $h=1$ part, and both weights in
\eqref{eq:Apoly} equal $\Gamma(\tfrac95)/\Gamma(\tfrac45)=\tfrac45$.  The
values at $t_{3,3}^{3}$ and $t_{3,3}^{4}$ solve the condition $\Ascr=0$ on the
diagonal of that monomial, with the lower values substituted into the terms
with $h\ge2$: at $t_{3,3}^{3}$ the term with $h=1$ is the single unknown
coefficient, of weight $1$, and at $t_{3,3}^{4}$ ($N=1$, balanced boundary
degrees $(2,7)$ and $(7,2)$) the minimal system sets the coefficient of
$(7,2)$ to zero and the condition determines that of $(2,7)$.
\end{remark}

\begin{remark}[The two-insertion values, and {\cite[Ex.~5.5]{GKTsurvey}}]
\label{rem:ex55}
Imposing \cite[Thm.~0.3]{GKT} directly on the oscillatory expansion of the
potential (the mechanism by which \cite[Thm.~0.2]{GKT} is proved, so that no
combinatorial convention enters) forces, at every two-insertion diagonal
$\{(a_{1},b_{1}),(a_{2},b_{2})\}$ with $a_{1}{+}a_{2}=r$, $b_{1}{+}b_{2}=s$,
$\mathbf d=0$, the relation \eqref{eq:seed}: the diagonal supports the boundary
degrees $(r,0)$ and $(0,s)$ only, $d(J,0)<0$ so the mirror theorem admits no term
there, and the linear and quadratic insertions contribute
$\nu_{(r,0)}/r+\nu_{(0,s)}/s$ and $1/(rs)$ respectively
(Lemma~\ref{lem:periods}).
At $r=s=4$ this gives $-\tfrac14$ for the sum considered in
\cite[Ex.~5.5]{GKTsurvey}, where $-\tfrac18$ is stated; the discrepancy traces to
the $h=2$ partition count in the evaluation of $\Ascr$ there, which
\cite[Notation~3.42 and Thm.~4.9]{GKT} fixes as the count of ordered partitions of the labeled
marking set with $1/h!$.  Notation~5.3 of \cite{GKTsurvey} writes the
partition sum over the multiset of insertions; read as ordered decompositions
of the multiset, the $h=2$ term is counted once and gives $-\tfrac18$, while
read over ordered partitions of a labeled realization it is counted twice and
gives $-\tfrac14$.  The same count is determined independently by
\cite[Cor.~5.11(3)]{GKT}: at $r=s=5$, $\Ascr(\{(1,1),(1,1)\},0)$ equals the closed
three-point value $\langle(\tau_{0}^{(1,1)})^{3}\rangle=1$, which the labeled
count reproduces and the halved count does not.  (The corresponding statement
in \cite[Cor.~0.6]{GKT} carries an extra sign $(-1)^{d(J,\mathbf d)-1}$.  The
diagonal of $t_{1,1}^2$ has $N=-1$, so the $t_{1,1}^2$-coefficient of
$\int_{\Xi_{2,2}}e^{W/\hbar}$ is $P_2(2)^2/(2\hbar^2)=1/(2\hbar^2)$; by
\cite[Thm.~0.3]{GKT} this equals $\hbar^{-2}\langle(\tau_{0}^{(1,1)})^{3}\rangle^{\mathrm{ext}}/2$,
which confirms the sign of \cite[Cor.~5.11(3)]{GKT}.) On
the marginal two-insertion diagonals the two
coefficients in \eqref{eq:seed} are moved (antisymmetrically) only by the torus
direction of Remark~\ref{rem:torus}; at $r=s$ the point of this
one-parameter set of values with equal coefficients has
$\nu_{(r,0)}=\nu_{(0,s)}=-\tfrac1{2r}$. The discrepancy concerns
only the worked example; no theorem of \cite{GKT} or \cite{GKTsurvey} is
affected.
\end{remark}

\section{Extremal systems and the boundary-ray factorization
\texorpdfstring{$\Dscr^{r,s}$}{D(r,s)}}
\label{sec:canonical}

\begin{proof}[Proof of Theorem~\ref{thm:canonical}]
Work in a primary quotient modulo $F^K$, or more generally on a finite
divisor-closed coefficient set as in Theorem~\ref{thm:presentation}(ii).
For existence and uniqueness of $\nu^{\min}$, order its monomials by total
degree. At each step, on each wall-bearing diagonal the span of the wall
vectors $v_{p}=\lambda_{p}\varepsilon_{p}-\mu_{p}\varepsilon_{p-1}$ ($\lambda_{p},\mu_{p}>0$) is complementary to
$\QQ\varepsilon_{0}$: a combination $\sum c_{p}v_{p}\in\QQ\varepsilon_{0}$ forces, reading the
$\varepsilon_{N},\varepsilon_{N-1},\dots$ components in turn, $c_{N}=\dots=c_{1}=0$. Hence there is a
unique group correction at this order making all $\nu_{p}$, $p\ge1$, vanish, and
it does not disturb lower orders; on marginal diagonals the torus direction plays
the role of the single wall term.  This triangular elimination is the argument of
\cite[proof of Thm.~4.27, Step~3]{GKT}, with the torus direction adjoined. The resulting system is realizable because the group
preserves $\Om(I)$ (Lemma~\ref{lem:enlarged}). Uniqueness follows from
freeness (Lemma~\ref{lem:enlarged}) plus the same triangularity: a nontrivial leading correction moves some
$\nu_{p\ge1}$. The construction for $\nu^{\max}$ is symmetric (complement of
$\QQ\varepsilon_{N}$). The factorization, in increasing order of slope, of the unique
$g^{\min\to\max}$ (Theorem~\ref{thm:presentation}(ii)) is then determined
without choices, and its ray factors are computed order by order from the
extremal systems. Uniqueness
makes these constructions compatible under restriction, giving the primary
completion and the coefficientwise descendent extension.
\end{proof}

\begin{example}[The first walls of $\Dscr^{5,5}$]
\label{ex:walls}
At order two in the deformation variables the three diagonals of
$x^{5}+y^{5}$ carrying positive-weight wall terms have balanced boundary degrees
$(0,6),(5,1)$; $(6,0),(1,5)$; and $(1,6),(6,1)$ respectively, with the
invariant relations and extremal values below (\S\ref{sec:cert}).  The two
marginal diagonals $u=t_{2,2}t_{3,3}$ and $u=t_{2,3}t_{3,2}$, with balanced
boundary degrees $(0,5)$ and $(5,0)$, carry the torus terms $\tfrac1{25}uX_{0,0}$
on the $(1,1)$-ray, by the two-insertion relation \eqref{eq:seed}:
{\small
\[
\renewcommand{\arraystretch}{1.35}
\begin{array}{lll}
\toprule
\text{monomial} & \text{relation} & \text{extremal values}\\
\midrule
t_{2,3}t_{3,3} &
\tfrac25\nu_{(0,6)}+\tfrac15\nu_{(5,1)}=-\tfrac2{25} &
\nu^{\min}_{(0,6)}=-\tfrac15,\ \nu^{\max}_{(5,1)}=-\tfrac25\\
t_{3,2}t_{3,3} & \text{the image under $x\leftrightarrow y$} & \text{by symmetry}\\
t_{3,3}^{2} &
\tfrac25\bigl(\nu_{(1,6)}+\nu_{(6,1)}\bigr)=-\tfrac4{25} &
\nu^{\min}_{(1,6)}=\nu^{\max}_{(6,1)}=-\tfrac25\\
\bottomrule
\end{array}
\]
}
Each relation is the condition $\Ascr=0$ of \eqref{eq:Apoly}.  For
$t_{2,3}t_{3,3}$ the $h=1$ weights are
$\Gamma(\tfrac75)/\Gamma(\tfrac25)=\tfrac25$ and
$\Gamma(\tfrac65)/\Gamma(\tfrac15)=\tfrac15$, and the $h=2$ term, the product
of the two one-marking normalizations, has weight
$\Gamma(\tfrac65)\Gamma(\tfrac75)/(\Gamma(\tfrac15)\Gamma(\tfrac25))=\tfrac2{25}$;
for $t_{3,3}^{2}$ the corresponding weights are $\tfrac25$, $\tfrac25$ and
$(\tfrac25)^2$.  In $\nu^{\min}$ only the coefficient of least $x$-degree
survives, and in $\nu^{\max}$ only that of greatest $x$-degree.  Solving
$W^{\nu^{\max}}-W^{\nu^{\min}}=c\,\cdot u\,X_{k_{1},k_{2}}(W_{0})$ on each
diagonal is an overdetermined system (two monomial components, one
coefficient); it is consistent because both systems satisfy the same linear
relation, whose solutions differ by multiples of the wall vector, and it
yields the ray factors of Theorem~\ref{thm:canonical}. The wall term
$t_{3,3}^{2}X_{1,1}$ is the first symmetric positive-weight primary wall term:
only the sum \eqref{eq:firstwallrel} is invariant,
the difference changes across the $(1,1)$-ray.
\end{example}

\section{Transport of the potential and realization by homotopies}
\label{sec:geometric}

The results above define the boundary-ray factorization
$\Dscr^{r,s}$ and show that its action, together with the period relations,
determines the invariants. What they do not provide is an
\emph{affine base}: an intrinsic domain in which the walls are rays and the
invariants are computed by broken lines, as for the Fano Landau--Ginzburg models
of \cite{Gross10,CPS}. This section builds as much of that picture as is a theorem
and isolates what is not.

The product $g^{\min\to\max}$ is the transport
along a path from one edge of $C_\partial$ to the other, whose two ends carry
$\nu^{\min}$ and $\nu^{\max}$; it is not the holonomy of a loop
(Proposition~\ref{thm:positionobstruction7}).

The section is organized as follows.  Section~\ref{ssec:sector} places the
walls; \S\ref{ssec:transport} proves the transport representation and its
ordered-chain inversion, and shows that the theory is strictly richer than the
part generated from $x^r+y^s$ by wall-crossing; \S\ref{ssec:realization} realizes every finite coefficient
quotient by concatenated homotopies and proves the obstruction that makes the
concatenation necessary; and \S\ref{ssec:missing} describes the punctured
quadrant (Proposition~\ref{thm:positionobstruction7}) and what it does not
supply (Remark~\ref{rem:interpretation}).

\subsection{The quadrant and its rays}
\label{ssec:sector}
\begin{definition}
\label{def:sector}
The \emph{boundary-degree lattice} is
$M_{\partial}=\ZZ e_{1}\oplus\ZZ e_{2}$, where $e_i$ records one boundary
insertion of type $\sigma_i$.  The \emph{boundary-degree quadrant} is its
effective cone $C_{\partial}=\RR_{\ge0}e_1+\RR_{\ge0}e_2$; its two \emph{edges}
are the rays of the two boundary types. A factor of $\Dscr^{r,s}$ is assigned
to the boundary ray $\RR_{\ge0}(k_{1}{+}1,k_{2}{+}1)$ of its critical graph,
whose direction is the Lie bidegree shifted by $(1,1)$
(Proposition~\ref{prop:shift}), and a chamber of a finite primary truncation is a connected
component of the complement of these rays. For the completed descendent
factorization, whose walls are dense (Proposition~\ref{prop:localfinite}), a
\emph{slope cut} is a direction together with a choice of side when the
direction itself carries a factor, that is, a point of
$\mathbb P(C_\partial)$ doubled at each wall slope; for transport it matters
only through the Dedekind cut it induces on the ordered set of wall slopes,
which in a finite truncation is a chamber, or a wall ray approached from one
side.  The \emph{$\sigma_1$-edge cut} and the \emph{$\sigma_2$-edge cut} are
the cuts at the edges of $C_\partial$ along the $\sigma_{1}$-axis and the
$\sigma_2$-axis, and the transport $\Psi(q)$ to a cut $q$ is the
slope-ordered product of the ray
factors between the $\sigma_1$-edge cut and $q$, defined on each finite
divisor-closed coefficient set and
then coefficientwise: on the coefficient of a fixed deformation monomial $u$,
every factor is projected to the finite divisor-closed quotient consisting of
the divisors of $u$, where only finitely many projected ray factors are
nontrivial.  Only the products $\Psi(q)\Psi(q')^{-1}$ enter below,
so the choice of the $\sigma_1$-edge cut as base point is a convention.  For
$\alpha,\alpha'\in\ZZ^{2}$ we write
$\det(\alpha,\alpha')=\alpha_{1}\alpha'_{2}-\alpha_{2}\alpha'_{1}$.
\end{definition}

Thus the quadrant, its lattice, and its edges come from the open FJRW boundary
data rather than from an auxiliary drawing. Only its projectivization
$\mathbb P(C_{\partial})\cong[0,\infty]$, the ordered slope space, is used by
transport. The placement is not a convention but forced
(Proposition~\ref{prop:placement}), and the same computation shows that the
$(1,1)$-shift of Proposition~\ref{prop:shift}, the removal of the
$\sigma_{12}$-insertion, is what makes bending along a transported monomial
monotone in slope (Theorem~\ref{thm:nobackscatter}).


\begin{proposition}[Forced placement]
\label{prop:placement}
The vector field $X_{k_{1},k_{2}}$ fixes exactly the monomials proportional to
its boundary-monomial direction $(k_{1}{+}1,k_{2}{+}1)$.  Thus the boundary
ray of a critical graph is the unique ray whose monomials its wall term
fixes, which is why the walls are placed there, and a monomial is fixed by the
walls on the ray of its own exponent.
\end{proposition}

\begin{proof}
$X_{k_{1},k_{2}}(x^{\alpha_1}y^{\alpha_2})=\bigl((k_{2}{+}1)\alpha_1
-(k_{1}{+}1)\alpha_2\bigr)x^{\alpha_1+k_{1}}y^{\alpha_2+k_{2}}$, and the
coefficient is the determinant of $\alpha=(\alpha_1,\alpha_2)$ against
$(k_{1}{+}1,k_{2}{+}1)$, vanishing precisely on the ray.
\end{proof}

If the two edges of the quadrant were identified, their ordered crossing would
define a loop with holonomy $g^{\min\to\max}\ne\mathrm{id}$
(Theorem~\ref{thm:canonical}); in $C_\partial$ the edges remain distinct
(Proposition~\ref{thm:positionobstruction7}).

\begin{proposition}[Interlacing and extreme support]
\label{prop:interlace}
On every diagonal the balanced monomials and the critical graphs interlace in
slope (writing $\mathrm{bal}_{p}$ for the balanced monomial of $\Gamma_{J,p}$),
$\mathrm{bal}_{N}<\Lambda_{N}<\dots<\Lambda_{1}<\mathrm{bal}_{0}$, and the
extremal systems of Theorem~\ref{thm:canonical} are the realizable systems at
the two extreme chambers of every finite primary truncation, with the support
of that theorem: near one edge only the monomial at the far end of each
diagonal survives.
\end{proposition}

\begin{proof}
Immediate from the formulas of Definition~\ref{def:diagonal}: the critical
graph $\Lambda_{p}$ has boundary degree between those of $\mathrm{bal}_{p}$ and
$\mathrm{bal}_{p-1}$ in slope; the support statement is the construction of
the extremal systems read through the interlacing.
\end{proof}



\subsection{The transport representation}
\label{ssec:transport}

\begin{definition}
\label{def:anchor}
The \emph{anchor cut} $q_m$ of a potential monomial $m$ is the one-sided cut
through its own direction, on the side toward the $(1,1)$-direction. If $m$ is
off every wall in a finite truncation, this is represented by the chamber
containing its direction, and if $m$ lies on a wall, by the chamber adjacent
to that wall on the side of the $(1,1)$-direction; for $m$ on the
$(1,1)$-direction either side gives the
same coefficient, and the exponent $(0,0)$ is anchored at the $\sigma_1$-edge
cut. Its
\emph{anchored value} is its coefficient at $q_m$. For a
candidate assignment of initial coefficients,
$\widetilde W(q):=\sum_{m}\Psi(q)\Psi(q_m)^{-1}(\mathrm{init}_{m}\,m)$ denotes the
resulting transport potential; the \emph{anchored-value assignment} sets each
$\mathrm{init}_m$ equal to the anchored value of $m$.
\end{definition}

Every bend along a transported monomial adds $\beta-(1,1)$ to the exponent,
$\beta$ the boundary degree of the wall term used, which lies on the crossed
ray.  As in the introduction, a \emph{bend chain} is a term of the
ordered bend expansion of Lemma~\ref{lem:bendchain} below: a sequence of wall
terms applied to a transported monomial along a monotone path in the slope
line.  Throughout, initial
monomial and target are taken off the crossed rays: an initial monomial on a
bend ray has vanishing bend coefficient, and an on-ray target uses the one-sided anchor cut toward
the $(1,1)$-direction, the convention of Definition~\ref{def:anchor}. Call a
bending contribution to a coefficient \emph{foreign} if it originates from an
initial coefficient other than the coefficient's own.

\begin{theorem}[Monotonicity of bends]
\label{thm:nobackscatter}
Along a path with $n_b$ bend factors, using wall terms of boundary degrees
$\beta_f$, from an initial monomial of exponent $\alpha'$ into a
monomial of exponent $\alpha$,
\[
n_b\det(\alpha,(1,1))
=\det(\alpha,\alpha')+\sum_{f}\det(\alpha,\beta_{f}).
\]
Consequently no bending contribution originating from a different initial
coefficient reaches a monomial's anchor cut from the far
side of the monomial's own direction, the side away from the
$(1,1)$-direction: at its anchor cut, a monomial below the $(1,1)$-direction
receives foreign bending only along paths descending in slope, one above it
only along paths ascending, and a monomial on the $(1,1)$-direction receives
none.
\end{theorem}

\begin{proof}
Summing the per-bend exponent motion gives
$\alpha+n_b(1,1)=\alpha'+\sum_{f}\beta_{f}$; the displayed identity is obtained
by applying $\det(\alpha,\cdot)$ to it. Transport is monotone, so every crossed ray lies strictly
between the directions of $\alpha'$ and $\alpha$. For a path from below, the
direction of $\alpha'$ and all crossed rays are strictly shallower than $\alpha$, so the right side is
negative, forcing $\det(\alpha,(1,1))<0$, i.e.\ $\alpha$ strictly above
the $(1,1)$-direction; from above it is positive, forcing $\alpha$ strictly below. A monomial on the
$(1,1)$-direction, where $\det(\alpha,(1,1))$ vanishes, excludes both.
\end{proof}

\begin{lemma}[Reduction]
\label{lem:reduction}
For the anchored-value assignment, $\widetilde W=W$ at every cut if and only if, for
every coefficient, the foreign bending contributions cancel at its anchor cut.
\end{lemma}

\begin{proof}
Ray factors act additively, so $W-\widetilde W$ is itself transported by them; its
coefficients on any divisibility-minimal nonzero monomial are
cut-independent (all bending contributions come from proper divisors, where
$W-\widetilde W$ vanishes), and evaluating each coefficient at its anchor cut gives zero
exactly under the stated cancellation.
\end{proof}

We call $x^{r}$ and $y^{s}$, the two monomials of $W_0$, the \emph{axis
monomials}.  For a twist $(a,b)$, the \emph{tower} of $(a,b)$ is the set
of singleton diagonals of $t_{a,b,d}$, $d\ge1$.

\begin{theorem}[Linear order and towers]
\label{thm:linearorder}
At linear order in the wall terms over a diagonal's monomial applied to the axis
monomials $x^{r},y^{s}$, the cancellation of
Lemma~\ref{lem:reduction} holds on every diagonal of every
$(r,s,J,\mathbf d)$: transporting into the cut at $\mathrm{bal}_{p}$ from below
crosses exactly the rays of the $\Lambda_{q}$, $q\ge p{+}1$, whose $x^{r}$-bends
land on $\mathrm{bal}_{q}$, never on the target coefficient, while from above
the $y^{s}$-bends at the crossed rays land on $\mathrm{bal}_{q-1}$, $q\le p$, again off
the target. For singleton
diagonals (one internal marking, any descendent) no contribution from another
diagonal exists,
the linearization is exact, and the transport representation holds with the
anchored-value assignment on every such tower, with wall coefficients the
chain $c_{q+1}=c_{q}\lambda_{q}/\mu_{q+1}$, with $c_1$ fixed by the condition
$\Ascr=(-1)^{d}$ of the diagonal, where
$\lambda_q,\mu_q$ are the coefficients of the wall vectors
$v_q=\lambda_q\varepsilon_q-\mu_q\varepsilon_{q-1}$ of \S\ref{sec:diagonal}.
\end{theorem}

\begin{proof}
The crossing sets and landing indices are those of
Proposition~\ref{prop:interlace} together with the landing identities: by
\eqref{eq:wallvector}, $X_{\Lambda_{q}}(x^{r})$ is a nonzero rational
multiple of the monomial $\mathrm{bal}_{q}$ and $X_{\Lambda_{q}}(y^{s})$ of
$\mathrm{bal}_{q-1}$. For a singleton
monomial no proper divisor exists, so the only contributions are the linear
ones.
\end{proof}

The interlacing pattern of Theorem~\ref{thm:linearorder} is the analogue for
this transport of the monotonicity of broken lines in \cite{CPS}.  Bends of
initial coefficients of divisors by the wall terms of other diagonals lie outside
the theorem; Example~\ref{ex:farside} shows that they can arrive.

Two computed examples below (Examples~\ref{ex:middleseeds}
and~\ref{ex:farside}) show what the initial coefficients are.

\begin{theorem}[Transport representation; $=$ Theorem~\ref{thm:transport}]
\label{thm:transport7}
With the anchor assignment of Definition~\ref{def:anchor}, the difference between
the potential's coefficient vector on a monomial and the transported contributions
of the initial coefficients of its proper divisors is independent of the cut.
Defining the \emph{initial system} by these constants, every cut potential is
the transport
\[
W(q)=\sum_{m}\Psi(q)\,\Psi(q_m)^{-1}(\mathrm{init}_{m}\,m)
\]
of the initial coefficients, and the initial system is the unique one with
this property.
\end{theorem}

\begin{proof}[Proof of Theorem~\ref{thm:transport7} ($=$ Theorem~\ref{thm:transport})]
Induct on the divisibility order of deformation monomials; the induction
starts at $W_{0}=x^{r}+y^{s}$, whose two coefficients are initial coefficients
with value $1$ in every cut. Fix a monomial $u=u_{J,\mathbf d}$ and assume the
initial coefficients are defined,
with the transport identity holding, on every proper divisor of $u$. Write
$D$ for the diagonal of $u$, so that $u_{D}=u$, and $E$ for the diagonal of a
proper divisor $u_{E}$ of $u$. Let
$W_{<u}(q)$ denote the $u$-component of the sum of the transported lower
initial coefficients.

\emph{Step 1: the change across a wall.}
Across any wall $\theta=\exp(\eta)$ the coefficients on $D$ change by
\[
\Delta_{\theta}\bigl(W|_{D}\bigr)
=\sum_{j\ge1}\frac1{j!}
\sum_{\substack{\eta_{1},\dots,\eta_{j}\in\operatorname{supp}(\eta)\\
u_{E}\,u_{\eta_{1}}\cdots\,u_{\eta_{j}}=u_{D}}}
\bigl[\eta_{1}\cdots \eta_{j}\bigl(W|_{E}\bigr)\bigr]_{D},
\]
a universal expression in $\eta$ and the coefficients of $W$ on
proper divisors $E$ of $D$ and on $W_{0}$: any other source overshoots
$u$, and the wall terms over $u$ act on $D$'s coefficients only through $W_{0}$
(their monomial is $u$ itself), so the change is exactly linear there.

\emph{Step 2: cancellation, and chamber-independence of the residual.}
The change of $W_{<u}$ across the same wall is the same universal expression
evaluated on the sum of transported lower initial coefficients, whose divisor
coefficients equal
$W$'s by the induction hypothesis and whose $W_{0}$-coefficients are the same
constants $1$. Hence the two changes cancel in the difference and
$V:=W(q)|_{D}-W_{<u}(q)$ is independent of $q$.

\emph{Step 3: the initial coefficient at $D$, and uniqueness.}
A $u$-term transports with unchanged
$u$-component (every wall action on it overshoots $u$), so taking as initial
coefficients $\mathrm{init}_m$ of the monomials $m$ on $D$ the components of
$V$, placed at the assigned anchor cuts, gives
$\widetilde W|_{D}=V+W_{<u}=W|_{D}$ at every cut (all sums are coefficientwise finite: only
finitely many initial coefficients of divisors and bend factors contribute to
a fixed monomial),
completing the induction. Uniqueness is
immediate: $V$ is determined by $W$ and the lower initial coefficients. Anchor
assignments affect only the higher-order terms of the transports, hence the
initial coefficients at higher orders, never the existence or uniqueness of the decomposition; the
$(1,1)$-side assignment is the one compatible with
Theorem~\ref{thm:nobackscatter}.
\end{proof}

\emph{Consequences.}  (a) On a coefficient receiving no foreign bending at its
anchor cut, the initial coefficient equals its anchored value.  (b) An extreme
coefficient of a wall-bearing diagonal that receives no foreign bending at the
cut where it vanishes in the corresponding extremal system, namely the
$\sigma_1$-edge cut for the coefficient of greatest $x$-degree and the
$\sigma_2$-edge cut for that of least $x$-degree, has initial coefficient
zero.

\begin{proof}[Proof of (a) and (b)]
Evaluate the transport formula of Theorem~\ref{thm:transport} at a cut $q$.
The term of the monomial $m$ itself contributes $\mathrm{init}_m$ to the
coefficient of $m$, because every wall term has a nonconstant deformation
coefficient and therefore moves a term only into strictly larger deformation
monomials; every other contribution is foreign.  At $q=q_m$ this gives (a), as
in Lemma~\ref{lem:reduction}.  At the $\sigma_1$-edge cut the cut potential is
$W^{\nu^{\min}}$, and at the $\sigma_2$-edge cut it is
$g^{\min\to\max}(W^{\nu^{\min}})=W^{\nu^{\max}}$; there the coefficient of
greatest, respectively least, $x$-degree of every wall-bearing diagonal
vanishes (Theorem~\ref{thm:canonical}; Proposition~\ref{prop:interlace}
identifies the extremal systems with the two extreme chambers), which gives
(b).  Theorem~\ref{thm:nobackscatter} restricts the direction from which foreign
bending can arrive: none reaches the anchor cut of a monomial from beyond its
own direction, and none reaches the anchor cut of a monomial on the
$(1,1)$-direction.
\end{proof}

\begin{lemma}[Ordered bend expansion]
\label{lem:bendchain}
Work in a finite divisor-closed coefficient quotient and write
$z^\alpha=x^{\alpha_1}y^{\alpha_2}$.  Suppose a transport path crosses, in chronological
order,
\[
 \Phi_j=\exp(\!\epsilon_j\eta_j),\qquad
 \eta_j=\sum_{f\in L_j}c_f u_f X_{k_f},
 \qquad \epsilon_j\in\{1,-1\}.
\]
For nonnegative integers $n_j$ and ordered choices
$f_{j,1},\ldots,f_{j,n_j}\in L_j$, read the decorations first by
increasing $j$ and put $\alpha^h=\alpha^{h-1}+k_{f_h}$, starting from $\alpha^0=\alpha$.
Then
\begin{equation}
\label{eq:bendchain}
 (\Phi_M\circ\cdots\circ\Phi_1)(u_Ez^\alpha)
 =\sum_\gamma
 \frac{\prod_h\epsilon_{j(h)}c_{f_h}
  \det(\alpha^{h-1},k_{f_h}+(1,1))}
 {\prod_j n_j!}\,
 u_E\prod_hu_{f_h}\,z^{\alpha^{|\gamma|}},
\end{equation}
the finite sum over the surviving decorated bend chains $\gamma$.
\end{lemma}

\begin{proof}
Expand each exponential in powers of $\eta_j$ and distribute each power over
its summands.  Successive applications use
\[
 X_k(z^\alpha)=\bigl((k_2+1)\alpha_1-(k_1+1)\alpha_2\bigr)z^{\alpha+k}
 =\det(\alpha,k+(1,1))z^{\alpha+k}.
\]
The exponential supplies $1/n_j!$, the path supplies the displayed factor
order, and nilpotence makes the sum finite.  Wall terms on one boundary ray
commute, so repeated bends may equivalently be grouped as multisets with the
same factorials.  The expansion is by chains: transport is linear
in its initial monomial.  Rooted Lie trees arise only after an additional
BCH expansion of the ray factors, which depends on choices.  \qedhere
\end{proof}

Write each monomial of the potential as $m=u_Dz^\alpha$, where $D$ is a
diagonal, $u_D$ its deformation monomial, and $\alpha$ an exponent on it, and
let $q_m$ be its anchor cut.  For a second monomial $m'=u_Ez^{\alpha'}$, define
the strictly divisibility-triangular transport kernel
\begin{equation}
\label{eq:transportkernel}
 \mathcal K_{mm'}=
 [u_Dz^\alpha]\,\Psi(q_m)\Psi(q_{m'})^{-1}(u_Ez^{\alpha'})
 \quad\text{when }E\mid D, E\ne D,
 \qquad \mathcal K_{mm'}=0\quad\text{otherwise}.
\end{equation}
Lemma~\ref{lem:bendchain} is an explicit formula for $\mathcal K_{mm'}$: retain the
chains whose deformation labels multiply $u_E$ to $u_D$ and whose exponent
increments carry $\alpha'$ to $\alpha$.

\begin{theorem}[Triangular inversion]
\label{thm:seedinversion}
On every finite divisor-closed coefficient set, let $\mathrm{anc}_m$ be the
anchored value of the monomial $m$ and $\mathrm{init}_m$ its initial
coefficient.  In vector form,
\begin{equation}
\label{eq:seedmatrix}
 \mathrm{anc}=(1+\mathcal K)\,\mathrm{init},\qquad
 \mathrm{init}=(1+\mathcal K)^{-1}\mathrm{anc}=\sum_{n\ge0}(-\mathcal K)^n\mathrm{anc}.
\end{equation}
In particular,
\begin{equation}
\label{eq:seedchain}
 \mathrm{init}_m=\mathrm{anc}_m+
 \sum_{n\ge1}(-1)^n
 \sum_{m_0<\cdots<m_n=m}
 \mathcal K_{m_nm_{n-1}}\cdots \mathcal K_{m_1m_0}\mathrm{anc}_{m_0},
\end{equation}
where $<$ includes strict divisibility of the deformation monomials and the
intermediate boundary exponents are summed.  Thus every initial coefficient is
a finite alternating sum of anchored values weighted by ordered
bend chains.
\end{theorem}

\begin{proof}
Evaluate Theorem~\ref{thm:transport} at the anchor cut of $m$.  The initial
coefficient of $m$ contributes $\mathrm{init}_m$, since a positive-degree wall
term acting on it overshoots its deformation monomial.  The initial coefficient
of a monomial $m'<m$ contributes $\mathcal K_{mm'}\mathrm{init}_{m'}$, and no
other initial coefficient contributes.  Hence
$\mathrm{anc}_m=\mathrm{init}_m+\sum_{m'<m}\mathcal K_{mm'}\mathrm{init}_{m'}$.
The kernel strictly increases
deformation monomials and is nilpotent in the finite quotient, so its geometric
series inverse is finite and gives
\eqref{eq:seedmatrix}--\eqref{eq:seedchain}.  \qedhere
\end{proof}

The following examples show that the initial coefficients contain
information not generated by the ray factors themselves.  Consider the part of
the potential obtained by transport through the ray factors from the axis
monomials $x^r,y^s$ alone, that is, the part generated from $x^r+y^s$ by
wall-crossing.  In the language of theta functions \cite{GHK15,GHKK}, the
observation below says
that the potential is not the analogue of
$\vartheta_{(r,0)}+\vartheta_{(0,s)}$: it requires further summands, with
coefficients $\mathrm{init}_m$.  The observation rests on
Example~\ref{ex:middleseeds}, which follows it.
\begin{observation}
\label{obs:richer}
The open theory is strictly richer than the part generated from $x^r+y^s$ by
wall-crossing. Bend chains generate
exactly the part of the potential reachable by transport from the initial
coefficients; the initial system is input that the
ray factors do not produce, and it is not zero. This has no analogue in the
smooth Fano case of \cite{Gross10,CPS}, where the potential is generated by
broken lines from its initial monomials. An example is the invariant $\tfrac15$ of
\eqref{eq:newinvariants}, on a diagonal without wall terms, which no ray action
can produce (Example~\ref{ex:middleseeds}).  With the anchor convention of
Definition~\ref{def:anchor}, nonzero initial coefficients also occur at
coefficients of wall-bearing diagonals, for instance the value $-\tfrac{25}{16}$ of a middle
balanced monomial in the same example.
\end{observation}

\begin{proof}
By Theorem~\ref{thm:transport7} and Lemma~\ref{lem:bendchain}, every cut
potential is the sum of the bend chains from the initial coefficients; by the
uniqueness there, the part generated from $x^r+y^s$ differs from the potential
once an initial coefficient other than those of $x^r,y^s$ is nonzero.  In
Example~\ref{ex:middleseeds} the coefficient $\tfrac15$ receives no bending,
so by Consequence~(a) after Theorem~\ref{thm:transport7} it is such an initial
coefficient, and so is $-\tfrac{25}{16}$, by Theorem~\ref{thm:linearorder}.
\end{proof}

\begin{example}[Initial coefficients not generated by wall-crossing]
\label{ex:middleseeds}
Two things this example shows, for $x^5+y^5$: bending cannot generate
everything, and, with the anchor convention of Definition~\ref{def:anchor},
initial coefficients are not confined to diagonals without wall terms. For
the first,
the invariant
$\langle(\tau_{0}^{(3,3)})^{3}\sigma_{1}^{4}\sigma_{2}^{4}\sigma_{12}\rangle
=\tfrac15$ of \eqref{eq:newinvariants} has no bending source ($t$-monomials
multiply along a bend chain, and the unique candidate wall term fixes the
initial monomial of $t_{3,3}$ exactly), so such coefficients are initial data.
For the second, take the diagonal of $t_{3,3,2}$ in the tower of the twist
$(3,3)$.  Here $N=2$, the balanced boundary degrees are $(3,13)$, $(8,8)$ and
$(13,3)$, and the wall terms lie on the rays of $(4,9)$ and $(9,4)$, with wall
vectors $v_1=45\varepsilon_1-20\varepsilon_0$ and
$v_2=20\varepsilon_2-45\varepsilon_1$ in the basis
$\varepsilon_0,\varepsilon_1,\varepsilon_2$ of coefficients, whose coordinates
we write $\nu_0,\nu_1,\nu_2$ (the coefficients $\nu_p$ of the proof of
Theorem~\ref{thm:effective}).  The singleton
condition $\Ascr=(-1)^2$ has only its $h=1$ part, with weights
$\Gamma(\tfrac45)\Gamma(\tfrac{14}5)/\Gamma(\tfrac45)^2=\tfrac{36}{25}$ at the
two ends and $\Gamma(\tfrac95)^2/\Gamma(\tfrac45)^2=\tfrac{16}{25}$ in the
middle:
\[
\tfrac{36}{25}\,\nu_0+\tfrac{16}{25}\,\nu_1+\tfrac{36}{25}\,\nu_2=1 .
\]
In increasing slope the transport crosses the ray of $(9,4)$ and then that of
$(4,9)$, so the coefficient vectors in the three chambers are
$\nu^{\min}=\nu_0\varepsilon_0$, $\nu^{\min}+c_2v_2$ and
$\nu^{\min}+c_2v_2+c_1v_1=\nu^{\max}$.  The $\Ascr$-condition for $\nu^{\min}$
and the vanishing of the $\varepsilon_0$- and $\varepsilon_1$-components of
$\nu^{\max}$ form the linear system
\[
\begin{pmatrix}\tfrac{36}{25}&0&0\\[2pt] 1&-20&0\\ 0&45&-45\end{pmatrix}
\begin{pmatrix}\nu_0\\ c_1\\ c_2\end{pmatrix}
=\begin{pmatrix}1\\ 0\\ 0\end{pmatrix},
\]
whose solution is the extreme value
$\nu_0=\langle\tau_{2}^{(3,3)}\sigma_{1}^{3}\sigma_{2}^{13}\sigma_{12}\rangle
=\tfrac{25}{36}$ and equal wall coefficients $c_{1}=c_{2}=\tfrac5{144}$.  The
middle chamber contains the direction $(1,1)$ of the middle balanced boundary
degree $(8,8)$, and there the coefficient vector is
$\nu^{\min}+c_2v_2=(\tfrac{25}{36},-\tfrac{25}{16},\tfrac{25}{36})$.  The
middle coefficient $-\mu_2c_2$ is nonzero because $c_2=c_1\lambda_1/\mu_2=c_1
=\nu_0/\mu_1\ne0$, and by Theorem~\ref{thm:linearorder} it is the initial
coefficient of the middle balanced monomial,
$\langle\tau_{2}^{(3,3)}\sigma_{1}^{8}\sigma_{2}^{8}\sigma_{12}
\rangle^{\mathrm{mid}}=-\tfrac{25}{16}$, where the superscript denotes the
value in the middle chamber.  The $\Ascr$-condition holds in all
three chambers, since both wall vectors are annihilated by the weights.
\end{example}

\begin{example}[An anchored value need not be an initial coefficient]
\label{ex:farside}
Theorem~\ref{thm:nobackscatter} excludes foreign contributions only from beyond
the target's own direction; contributions from the $(1,1)$-side may remain at
the anchor cut.  Composite diagonals (those of monomials with at least two
internal markings) with descendents exhibit this behavior.  For $x^5+y^5$ we
consider the composite diagonals of $J=\{(0,0,2),(1,1,0)\}$, whose wall
term (the one over the composite monomial itself) lies on the $(1,1)$-ray, and
of $J=\{(0,0,2),(2,1,0)\}$, the smallest example whose wall term lies off the
$(1,1)$-ray, on the ray of $(3,2)$.  For each of them we solve
the $\Ascr$-conditions on the diagonals of all divisors of the composite
monomial and transport the resulting cut potentials across every ray carrying
a wall term of these finitely many diagonals: the diagonal of $t_{0,0,2}$ in
the tower of the twist $(0,0)$ carries
nonzero wall coefficients $\tfrac56$ on both of its rays, those of $(1,6)$ and
$(6,1)$; the maximal system forces the coefficient $1$ of the wall term over
each composite monomial; and the $\Ascr$-conditions hold in every chamber
(these exact calculations are listed in Section~\ref{sec:cert}).
In the second
case the balanced boundary degrees are $(2,6)$ and $(7,1)$; the coefficient of
$(7,1)$ vanishes in $\nu^{\min}$, so its anchored value, read below the ray of
$(6,1)$, is $0$.  The monomial $t_{2,1}x^2y$ has initial coefficient $1$ and
is anchored in the chamber between the rays of $(6,1)$ and $(3,2)$, so its
transport to that anchor cut crosses the ray of $(6,1)$ negatively, where the
logarithm of the ray factor is $\tfrac56\,t_{0,0,2}X_{5,0}$.  By \eqref{eq:bendchain} it
gives the foreign contribution
\[
-\tfrac56\,\det\bigl((2,1),(6,1)\bigr)\,t_{0,0,2}t_{2,1}x^{7}y
=\tfrac{10}{3}\,t_{0,0,2}t_{2,1}x^{7}y,
\]
the only one at that cut.  Hence the anchored value $0$, the value in the
minimal system, is not the initial coefficient, which is
$0-\tfrac{10}{3}=-\tfrac{10}{3}$.
\end{example}


\subsection{Realization by concatenated homotopies}
\label{ssec:realization}

In this subsection and the next, a \emph{leg} is one normal-form homotopy
(\S\ref{sec:presentation}) joining two consecutive intermediate canonical
families, that is, canonical families realizing consecutive members of a
sequence of realizable systems; the homotopies constructed below are
concatenations of legs.

\begin{proposition}[Descendent rays accumulate; primary rays are locally finite]
\label{prop:localfinite}
Fix $r,s\ge2$ and an interior twist $(a,b)$. Let $d$ range over $\ZZ_{\ge1}$, so that $t_{a,b,d}$ runs over deformation
parameters of total degree one. The rays of the associated critical graphs are
exactly
\[
R(a,b)=\bigl\{(a+j_1r+1,\ b+j_2s+1)\ :\ j_1,j_2\in\ZZ_{\ge0}\bigr\},
\qquad d=j_1+j_2+1,
\]
their coefficients in the ray factors of $\Dscr^{r,s}$ are all nonzero, and
their slopes are dense
in $(0,\infty)$. Hence the walls of $\Dscr^{r,s}$
accumulate in every open subcone of $C_\partial$ already at total degree one
in the deformation parameters, for every $(r,s)$, the simple singularities
included; those carry no primary wall at all (Theorem~\ref{thm:trichotomy}). By contrast, write
$\Dscr^{r,s}_{0}$ for the \emph{primary boundary-ray factorization}, the factors with
$\mathbf d=\mathbf 0$, and $\mathfrak m$ for the ideal generated by the
deformation parameters. Then $\Dscr^{r,s}_{0}$ has finitely many walls modulo
$\mathfrak m^{n}$ for every $n$, and its chambers are open angular sectors.
\end{proposition}

\begin{proof}
For a singleton $J=\{(a,b)\}$ one has $\ell_{1}=\ell_{2}=0$ and $N=d$, so by
Definition~\ref{def:diagonal} the critical graphs $\Lambda_{J,p}$, $1\le p\le d$,
have boundary degrees $(a+(p-1)r+1,\ b+(d-p)s+1)$; putting $j_1=p-1$,
$j_2=d-p$ gives $j_1+j_2=d-1$ and the stated set, every element of which is
attained by taking $d=j_1+j_2+1$. For density, fix $\xi>0$ and let $j_1$
be large enough that $\xi(a+j_1r+1)>1+b$; with
$j_2=\lceil (\xi(a+j_1r+1)-1-b)/s\rceil$ one has
$j_2 s\in[\,\xi(a+j_1r+1)-1-b,\ \xi(a+j_1r+1)-1-b+s)$, whence
\[
\Bigl|\frac{b+j_2 s+1}{a+j_1r+1}-\xi\Bigr|\ <\ \frac{s}{a+j_1r+1}\
\xrightarrow[\ j_1\to\infty\ ]{}\ 0 .
\]
For the coefficients: on a singleton diagonal the chamber-index
condition is $\Ascr(J,\mathbf d,\nu)=(-1)^{d}\ne0$, whose $h=1$ part is
$\sum_{p}\omega_{p}\nu_{p}$ with all $\omega_{p}>0$, so $\nu^{\max}_{N}\ne0$;
solving $W^{\nu^{\max}}-W^{\nu^{\min}}=\sum_{p}c_{p}v_{p}$ with
$v_{p}=\lambda_{p}\varepsilon_{p}-\mu_{p}\varepsilon_{p-1}$, $\lambda_{p},\mu_{p}>0$, gives
$c_{N}=\nu^{\max}_{N}/\lambda_{N}\ne0$ and $c_{p}=c_{p+1}\mu_{p+1}/\lambda_{p}\ne0$ for every
$p$, and no cancellation between distinct rays is possible at total degree one,
where the monomials are distinct.
For the primary statement, modulo $\mathfrak m^{n}$ only monomials
$u_{J,\mathbf 0}$ with $|J|<n$ occur, and for each such $|J|$ the twists range
over the finite range of narrow residues $0\le a\le r-2$, $0\le b\le s-2$, so
only finitely many diagonals, hence
finitely many walls, contribute; finitely many rays through the origin cut the
quadrant into open angular sectors.
\end{proof}

\begin{theorem}[Realization by homotopies in a finite coefficient quotient;
  $=$ Theorem~\ref{thm:georealization}]
\label{conj:geometric}
Let $S$ be a finite divisor-closed set of deformation monomials containing
$1$, and choose a finite labeled set $I$ with at least $n_{a,b}(S)$ labels of
each twist occurring in $S$, with $n_{a,b}(S)$ as in
Theorem~\ref{thm:georealization}.  The quotient factorization, the reduction to
$S$ of the factorization $\Dscr^{r,s}$, has finitely many walls, with rays
$\rho_1,\dots,\rho_L$ listed in increasing slope.  There is a homotopy
$(\mathfrak s_{T})_{T\in[0,1]}$ of boundary conditions, that is, a family of
simple canonical homotopies in the sense of \cite[Defs.~3.36--3.37]{GKT},
obtained by concatenating normal-form homotopies, from a family of canonical multisections realizing $\nu_I^{\min}$ to
one realizing $\nu_I^{\max}$ (the restrictions of the minimal and maximal
systems to $I$); every time slice is a family of canonical multisections
\cite[Def.~3.22]{GKT}.  After reduction to the quotient supported on $S$ every
nonzero jump is supported on a single boundary ray, the rays carrying nonzero
reduced jumps occur in increasing slope as $T$ increases, and the total
reduced jump on $\rho_{i}$
is the ray factor $\theta_{\rho_{i}}$ of
Theorem~\ref{thm:canonical} modulo that quotient.
\end{theorem}

\begin{proof}
\emph{Step 1: the labeled model, and symmetrization.}
Let $\mathfrak q_S$ be the monomial ideal spanned by monomials outside $S$;
it is an ideal because the complement of a divisor-closed set is upward
closed.  Write a bar for reduction modulo $\mathfrak q_S$. Let
$A_I=\QQ[u_{i,j}:i\in I,\ j\ge0]/(u_{i,j}u_{i,j'})$ be the labeled
coefficient ring of \cite[Def.~4.6]{GKT}, square-zero in each label; write
$u_{A,\mathbf d}=\prod_{i\in A}u_{i,d_i}$ and $u_{i}=u_{i,0}$. The monomials
$u_{J,\mathbf d}$ (indexed by a twist multiset, in the deformation
parameters) and $u_{A,\mathbf d}$ (indexed by a label subset, in $A_{I}$)
live in different rings.
Symmetrization sends, for every descendent level,
\[
t_{a,b,d}\longmapsto
\sum_{\substack{i\in I\\\operatorname{tw}(i)=(a,b)}}u_{i,d},
\]
where $\operatorname{tw}(i)$ denotes the twist of the label $i$.
Let $\mathfrak q_{I,S}$ be the labeled monomial ideal whose monomials have
symmetric type outside $S$.  Lemma~4.14 of \cite{GKT} says that symmetrization induces
an injection
\[
 \QQ[t_{a,b,d}]/\mathfrak q_S\hookrightarrow A_I/\mathfrak q_{I,S}.
\]
Indeed, a monomial in $S$ maps to the nonempty sum over injective assignments
of its occurrences to labels of the required twists; the saturation bound
$n_{a,b}(S)$ makes the sum nonempty, and distinct symmetric monomials have
disjoint labeled supports. The induced
map on the enlarged Lie algebra preserves boundary rays, including the
marginal $(1,1)$-ray. The action after reduction is still free: the
minimal-support argument of Lemma~\ref{lem:enlarged} detects every nonzero
class supported in $S$ on $W_0$.
Thus the $\mathfrak q_{I,S}$-reductions of $\nu_I^{\min},\nu_I^{\max}$ and of the labeled
boundary-ray factorization identify with their universal symmetric counterparts. Below we
suppress the labeled-lift symbol in the action of a symmetric factor.

\emph{Step 2: the ray factors are symmetric.}
Twist-preserving permutations fix the two symmetric extremal systems, hence by
freeness fix $g^{\min\to\max}$; they also preserve every boundary-ray
subspace. Uniqueness of the factorization in
Theorem~\ref{thm:presentation}(ii), whose filtered proof applies verbatim
over the labeled coefficient ring, therefore makes each
$\theta_{\rho_i}$ symmetric.
Write the boundary-ray factorization, in the composition convention of
\S\ref{ssec:results}, as
\[
g^{\min\to\max}
=\theta_{\rho_L}\circ\cdots\circ\theta_{\rho_1}
\pmod{\mathfrak q_{I,S}},
\qquad \operatorname{slope}(\rho_1)<\cdots<
\operatorname{slope}(\rho_L).
\]
\emph{Step 3: one leg per ray.}
For each quotient factor choose the symmetric full-group representative
$\widehat\theta_{\rho_i}=\exp(\widehat\eta_{\rho_i})$ whose logarithm is the
finite sum of its terms supported in $S$ and has no terms outside $S$.
Starting with a transverse canonical family realizing $\nu_I^{\min}$, that
is, a family of canonical multisections \cite[Def.~3.22]{GKT} transverse to
the zero section, as required of the endpoints in Lemma~\ref{lem:gktextension},
construct one leg for each factor. If the current system is $\nu$, then
$\widehat\theta_{\rho_i}\nu$ is again an element of $\Om(I)$. The
realizability statement $\Om(I)=\Upsilon_{r,s}(I)$ of \cite[Cor.~5.14]{GKT},
read in the enlarged sense of Lemma~\ref{lem:enlarged} on the marginal ray,
supplies a canonical endpoint family realizing it, and the constructive
realization in \cite[Thm.~5.13, Lem.~5.15]{GKT} may be chosen transverse.
Only the systems, not the families, need to be symmetric. Join the two
transverse endpoint families by a normal-form homotopy as in
Lemma~\ref{lem:gktextension}.

\emph{Step 4: every jump of the leg lies on one ray.}
We claim that every jump internal to this leg is supported on $\rho_i$.
For the actual chronological order of its subset times, the jump formula
(\S\ref{sec:presentation}) writes the total labeled jump as
\[
\prod_{\varnothing\ne A\subseteq I}^{\mathrm{time}}\exp(\eta_A),
\]
where $\eta_A$, the logarithm of the jump at $T_A$, has exact label support
$A$ (it is the sum over descendent vectors of terms $u_{A,\mathbf d}X$). In the labeled coefficient ring,
$u_{A,\mathbf d}u_{B,\mathbf e}=0$ when $A\cap B\ne\varnothing$.
For any fixed order this exact-support factorization is unique. Indeed,
a nonzero bracket of disjoint supports has support $A\cup B$; hence
\[
\left.\log\!\left(\prod_A\exp(\eta_A)\right)\right|_A
=\eta_A+P_A(\eta_B:B\subsetneq A),
\]
and induction on $\lvert A\rvert$ recovers every $\eta_A$.

The total jump and the labeled lift of $\widehat\theta_{\rho_i}$ carry the same initial
endpoint potential to the same final one. Freeness of the enlarged labeled
action (Lemma~\ref{lem:enlarged}) makes them equal. Now induct again on
$\lvert A\rvert$. The exact-$A$ component of the logarithm of the target is
supported on $\rho_i$; by induction all $\eta_B$ for $B\subsetneq A$ are supported
there as well. Every Lie word in $P_A$ therefore vanishes because the
boundary-ray subalgebra is abelian, as proved in
Theorem~\ref{thm:presentation}(ii). Thus $\eta_A$ equals the target's exact-$A$
component and is supported on $\rho_i$. This proves the claim, including for
the torus terms of the enlarged algebra.

\emph{Step 5: the final leg.}
After the legs of Step~3, the resulting full system $\nu^*$ is congruent to
$\nu_I^{\max}$ in the stated quotient. Let $g_{*}$ be the transition element, before
reduction, from $\nu^*$ to $\nu_I^{\max}$. Quotient freeness gives $\bar g_{*}=\mathrm{id}$. Join
the current endpoint family to a transverse canonical family realizing $\nu_I^{\max}$ by
one further leg, chosen as above. Its total jump equals $g_{*}$ by full
freeness. Applying the exact-support induction above \emph{after reduction} to
its time-ordered factorization shows successively that every one of its
internal jump factors has trivial class: this final leg is entirely
zero modulo $\mathfrak q_{I,S}$.

\emph{Step 6: concatenation.}
Give all legs constant collars at their endpoints, reparametrize them to
successive subintervals of $[0,1]$, and concatenate. Canonicity
\cite[Def.~3.22]{GKT}, strong positivity (the positivity conditions of
\cite[Def.~3.18]{GKT} at contracted boundary nodes and near the positive
boundary strata) and simplicity \cite[Defs.~3.36--3.37]{GKT} survive
this operation: use identical collars and
time reparametrizations simultaneously on every graph and component, which
preserves the boundary-factorization identity \cite[(3.12)]{GKT} and the
component-independence defining simplicity; the conditions (1)--(3) of
Lemma~\ref{lem:gktextension} hold in each leg, with leg-dependent times, and
no transversality is asserted for the components on which that lemma imposes
none; strong positivity is slicewise.
The resulting homotopy has the exact endpoints and the asserted reduced order
and total jumps. It is not one global normal-form homotopy: the same
marking subset may jump in different ray legs. That distinction is precisely
what evades Proposition~\ref{prop:obstruction}.

\end{proof}

\begin{remark}[Finite versus completed realization]
Theorem~\ref{conj:geometric} is a theorem for each fixed finite quotient.  The
ray factors are algebraically compatible when the finite divisor-closed
coefficient set grows,
but the proof does not choose their endpoint families and homotopies
compatibly.  Such a choice would require relative lifting theorems for the
spaces of canonical families and simple canonical homotopies.  We therefore
do not assert a single completed geometric homotopy for the dense descendent
factorization.  For $x^4+y^4$ and finite systems of graph data closed under
the assembling relation of \cite{GKT}, compatible choices are constructed in
\cite[Theorem~5.7 and Appendix~B]{annular}.  Moreover, the final leg is zero
in the stated quotient but may have jumps on several rays outside it.
\end{remark}

A single normal-form homotopy already suffices below the first degree at which
one diagonal carries two wall terms.

\begin{proposition}[Exact first two-wall order; realization by one
normal-form homotopy]
\label{prop:parabolicgeom}
For $n\ge1$, the maximum of the numerical diagonal count
$N(J,\mathbf0)$ over primary multisets with $\lvert J\rvert=n$ is
\[
N_{\max}(r,s;n)
=n+1-\left\lceil\frac{2n}{r}\right\rceil
-\left\lceil\frac{2n}{s}\right\rceil ,
\]
attained by $J=\{(r-2,s-2)^n\}$. Consequently the least primary degree
admitting a diagonal with two wall terms is
\[
N_2(r,s)=
\min\left\{n\ge3:
\left\lceil\frac{2n}{r}\right\rceil+
\left\lceil\frac{2n}{s}\right\rceil\le n-1\right\},
\]
where the minimum of the empty set is $\infty$.
\begin{enumerate}
\item[(i)] $N_{2}=\infty$ if and only if $\chat_{W}\le1$, that is, if and only if
$W$ is simple or parabolic; when $\chat_W>1$,
$N_{2}\ge\lceil 1/(\chat_{W}-1)\rceil$.
\item[(ii)] When $N_2<\infty$, a single normal-form
homotopy realizes the primary boundary-ray factorization modulo $\mathfrak m^{N_{2}}$. When
$N_2=\infty$, such a homotopy exists at every finite order, realizing the
completed primary boundary-ray factorization as the inverse system of these finite-order
statements; no compatibility of the homotopies themselves is asserted.
\end{enumerate}
\end{proposition}

\begin{proof}
Write $\sum_j a_j=r(J)+r\ell_1$ and
$\sum_j b_j=s(J)+s\ell_2$. Since $a_j\le r-2$,
\[
\ell_1\le\left\lfloor\frac{n(r-2)}r\right\rfloor
=n-\left\lceil\frac{2n}{r}\right\rceil,
\]
and similarly for $\ell_2$; both bounds are attained simultaneously by the
maximal twist. Substitution in
$N=\ell_1+\ell_2-n+1$ proves the formula for $N_{\max}$ and hence the displayed
formula for $N_2$. In particular $n=1,2$ never gives $N\ge2$.

The coarser bound from the proof of Theorem~\ref{thm:trichotomy} is
$\ell_{1}+\ell_{2}\le n\chat_{W}$, hence
\begin{equation}
\label{eq:Nbound}
N(J,\mathbf 0)\;\le\;1+n(\chat_{W}-1),
\end{equation}
which is at most $1$ for every $n$ when $\chat_{W}\le1$, and, when
$\chat_W>1$, forces $n\ge1/(\chat_{W}-1)$ whenever $N\ge2$; this is (i)
apart from the converse.
Conversely, if $\chat_{W}>1$, the exact maximal-twist count satisfies
$N_{\max}>n(\chat_{W}-1)-1$, which is at least $2$ for all sufficiently large
$n$, so $N_{2}<\infty$.

For (ii), put $K=N_2$ when $N_2<\infty$ and otherwise fix an arbitrary finite
$K$. Work modulo $\mathfrak q_K=(\mathfrak m^K,\text{descendent parameters})$ and choose
at least $K-1$ labels of every primary twist, so the symmetrization map used in
Theorem~\ref{conj:geometric} is injective in the degrees retained by the quotient. Then every
retained diagonal has $N\le1$, and a marking subset $A\subseteq I$ determines
the single primary diagonal $(A,\mathbf0)$. Choose transverse symmetric families
realizing $\nu_I^{\min}$ and $\nu_I^{\max}$ and take the one normal-form
homotopy between them. Its jumps occur at the subset times $T_A$. By the jump
formula of \cite[proof of Thm.~5.9]{GKT} (\S\ref{sec:presentation}), the
reduction of the jump at
$T_A$ is a sum over the critical graphs with marking set $A$: it is zero if
$|A|\ge K$ or no retained wall term exists, and otherwise is supported on the
unique wall term of that diagonal.

Order the times of all subsets whose retained primary diagonal carries its
unique wall term by the slope of that term's ray; place the remaining subset
times arbitrarily. A resulting zero jump is harmless. This prescription is legitimate:
\cite[Lem.~3.40]{GKT} constrains the times only to be distinct in $(0,1)$, and its
construction takes them as given: the zeros are moved onto a prescribed
$T_{I(\Gamma)}$ by a diffeomorphism preserving the moduli factor, and the
finitely many non-transverse times are kept off the $T_{A}$ by re-extending
transversally at each prescribed time. Visible jumps sharing a ray commute by the abelianness proved in
Theorem~\ref{thm:presentation}(ii). Full freeness identifies the total jump
with $g^{\min\to\max}$; after reduction, its grouping by ray is therefore the
unique slope-ordered factorization
$\prod^{\rightarrow}_{\rho}\theta_{\rho}$. Symmetry and injectivity descend
this labeled statement to the universal primary quotient. Thus the total
reduced jump on $\rho_i$ is $\theta_{\rho_i}$.
\end{proof}

\begin{remark}
\label{rem:orderinggap}
The ordering step above is the only place the freedom in the times of
\cite[Lem.~3.40]{GKT} is used; the times $T_A$ there may be prescribed
arbitrarily as distinct points of $(0,1)$. What the ordering does
\emph{not} do is separate the wall terms of a single diagonal: those all carry the one
marking set $J$, so they share the one time $T_{J}$ whatever order is chosen, which
is Proposition~\ref{prop:obstruction}.
\end{remark}

At the first two-wall degree, one global normal-form homotopy provably
does not suffice. This explains why the concatenation of legs in
Theorem~\ref{conj:geometric} is necessary.

\begin{proposition}[One normal-form homotopy cannot separate a diagonal's wall terms]
\label{prop:obstruction}
Let $(J,\mathbf 0)$ be a primary diagonal with $N\ge2$ and let $H^{*}$ be a family
of simple canonical homotopies with times $(T_{A})_{A\subseteq I}$ as in
\cite[Lem.~3.40]{GKT}. Then every possible contribution from the diagonal's
$N$ wall terms is concentrated at the single time $T_{J}$, with wall term $p$
actually contributing if and only if $\Delta(\Lambda_{J,p})\ne0$. The primary part of
the jump there, modulo the descendent parameters, is supported on one ray if
and only if at most one of these zero counts is nonzero.
In particular no single
such normal-form homotopy gives the ray-by-ray realization of
Theorem~\ref{conj:geometric} once two of the diagonal's wall terms actually
contribute.
\end{proposition}

\begin{proof}
The $N$ critical graphs $\Lambda_{J,1},\dots,\Lambda_{J,N}$ all carry the internal
marking set $J$, so by Lemma~\ref{lem:gktextension}(3)(a) their zeros occur at the one
time $T_{J}$. By the jump formula of \cite[proof of Thm.~5.9]{GKT}
(\S\ref{sec:presentation}), the logarithm of the jump there is
\[
\sum_{p=1}^{N}(-1)^{|J|+k_{2}(\Lambda_{J,p})}\,\Delta(\Lambda_{J,p})\,
u_{J,\mathbf0}X_{\Lambda_{J,p}},
\]
with $\beta(\Lambda_{J,p})=k(\Lambda_{J,p})+(1,1)$ the boundary degree of
$\Lambda_{J,p}$.  Its boundary rays are pairwise distinct: along $p$ the first
boundary degree increases by $r$ and the second decreases by $s$, both by at least
$2$ from values at least $1$, so the slope strictly decreases; and for $N\ge2$ none
of them is the torus direction. The displayed wall terms lie in distinct
bidegree summands, so no cancellation is possible. Its logarithm is therefore
supported on one ray exactly when at most one coefficient is nonzero.
\end{proof}

The concatenation of legs in Theorem~\ref{conj:geometric} is not a single
normal-form homotopy: it applies \cite[Lem.~3.40]{GKT} separately
to each boundary-ray factor and may therefore reuse the same marking subset at
several leg times. Proposition~\ref{prop:parabolicgeom} succeeds with a
single leg because when $N\le1$ the at-most-one condition on
$\Delta(\Lambda_{J,p})$ holds for free.

\begin{example}[The first two-ray diagonal of $x^5+y^5$]
\label{ex:tworay}
For $x^{5}+y^{5}$ the order $N_{2}=5$ of
Proposition~\ref{prop:parabolicgeom} is attained by $J=\{(3,3)^{5}\}$,
$\mathbf d=\mathbf 0$: here $m(J,\mathbf 0)=50$, the Euler weight is $25$, and
$N=2$, with balanced boundary degrees $(0,10),(5,5),(10,0)$ and critical graphs of
boundary degrees $(1,6)$ and $(6,1)$: two wall terms on two \emph{distinct} rays,
the first such diagonal for this Fermat pair. Their wall vectors are
$v_{1}=30\varepsilon_{1}-5\varepsilon_{0}$ and $v_{2}=5\varepsilon_{2}-30\varepsilon_{1}$. Since $|J|\ge2$ and $N\ge0$ (so
$d(J,\mathbf 0)=-N-1=-3<0$), the
chamber-index condition is $\Ascr(J,\mathbf 0,\nu)=0$, which reads
\[
\tfrac{6}{25}\,\nu_{(0,10)}+\tfrac1{25}\,\nu_{(5,5)}+\tfrac{6}{25}\,\nu_{(10,0)}
\;=\;\tfrac{204}{15625}:
\]
three coefficients against two torsor directions and one condition, as
Theorem~\ref{thm:effective} requires. In the minimal system of
Theorem~\ref{thm:canonical} this gives
\[
\bigl\langle(\tau_{0}^{(3,3)})^{5}\,\sigma_{2}^{10}\sigma_{12}\bigr\rangle^{\min}
=\tfrac{34}{625},
\]
with $\bigl\langle(\tau_{0}^{(3,3)})^{5}\sigma_{1}^{10}\sigma_{12}\bigr\rangle^{\max}$
the same by symmetry. This is the lowest primary degree
for $x^5+y^5$
at which the jump of a single normal-form homotopy would have to carry two rays,
hence the first primary degree at which the realization of
Theorem~\ref{conj:geometric} differs from one global normal-form homotopy
(Corollary~\ref{cor:nosimple}); every other primary diagonal computed
in \S\ref{sec:canonical} and in \S\ref{sec:geometric} has $N\le1$, which is
why the difficulty does not surface there.

The hypothesis of Proposition~\ref{prop:obstruction} is met here. Take $I$ to carry
exactly five markings of twist $(3,3)$, so that this monomial is realized by a single
marking set and the wall terms above it share one time. Writing
$\nu^{\max}-\nu^{\min}=c_{1}v_{1}+c_{2}v_{2}$ and reading the extreme rows gives
$c_{1}=c_{2}=\tfrac{34}{3125}$, the middle row being consistent; more generally
$c_{1}=\nu_{0}/\mu_{1}$ and $c_{p+1}=\lambda_{p}c_{p}/\mu_{p+1}$, where $\nu_0$ is
the $\varepsilon_0$-coefficient of $\nu^{\min}$ (here $\nu^{\min}_{(0,10)}$), so
every $c_{p}$ is nonzero as soon as $\nu_{0}$ is, exactly as in the proof of
Proposition~\ref{prop:localfinite}.
No lower-order term contributes to either coefficient: the wall terms over the
powers of $t_{3,3}$ sit at
$|J|=2,4,5$ and no multiset drawn from the lower orders $|J|=2,4$ sums to $5$, while a wall term acting on a lower
coefficient of the potential contributes nothing here, its direction being parallel to
that coefficient's and hence annihilating it (Proposition~\ref{prop:placement}). So both
zero counts are nonzero and both wall terms contribute to the jump at the one time.

The same computation shows the difficulty is not peculiar to the hyperbolic regime. A
single insertion of maximal twist with two descendents has $N=2$ for every $(r,s)$, with
both coefficients again nonzero, so the obstruction appears in every model once
descendents are admitted, the simple ones included, which carry no primary wall at
all. What is special about $x^{5}+y^{5}$ is that it is the first symmetric
hyperbolic Fermat pair and this is its first primary occurrence.
\end{example}

\begin{corollary}
\label{cor:nosimple}
For $x^{5}+y^{5}$ modulo $\mathfrak m^{6}$, and $I$ carrying at least five
labels of twist $(3,3)$, no normal-form homotopy supplied by
\cite[Lem.~3.40]{GKT} between families realizing the minimal and maximal
systems $\nu^{\min}$ and $\nu^{\max}$ has all of its jumps supported on single
rays.
\end{corollary}

\begin{proof}
By Example~\ref{ex:tworay} both wall terms of the diagonal $J=\{(3,3)^{5}\}$ contribute,
and by Proposition~\ref{prop:obstruction} they share the single time $T_{J}$, so the
jump there meets two distinct rays.
\end{proof}

\begin{remark}
\label{rem:thresholdgeom}
For simple $W$ the \emph{primary} statement carries no walls at all: $N\le0$
throughout by Theorem~\ref{thm:trichotomy}, so $\Dscr^{r,s}_{0}$ is empty and
the quadrant $C_\partial$ is a single chamber. By
Proposition~\ref{prop:localfinite}, though, the descendent walls of the same
model are dense in the same quadrant. The parabolic pairs are the first
case with primary content, and there the unique
ray is $(1,1)$, carrying the marginal torus direction; the change it effects is
the one exhibited in \cite[Ex.~4.9]{GKTsurvey}. In the hyperbolic regime diagonals
with $N\ge2$ occur and the jump of a single normal-form homotopy can then carry
several rays at once.
The exact first degree is given by Proposition~\ref{prop:parabolicgeom}.
Up to the symmetry $(r,s)\leftrightarrow(s,r)$, its complete hyperbolic
classification is
\[
\renewcommand{\arraystretch}{1.15}
\begin{array}{ccc}
\toprule
r&s&N_2(r,s)\\
\midrule
3&7&21\\
3&8&12\\
3&9\le s\le11&9\\
3&s\ge12&6\\
4&5&10\\
4&6\le s\le7&6\\
4&s\ge8&4\\
5&5\le s\le7&5\\
5&s\ge8&4\\
r\ge6&s\ge r&3\\
\bottomrule
\end{array}
\]
Thus $N_2\le21$ for hyperbolic integral Fermat exponents, with the maximum attained only
at $(3,7)$ up to symmetry. The coarser lower bound in
Proposition~\ref{prop:parabolicgeom} is useful for a fixed pair but does not
produce an unbounded sequence approaching the threshold: the exponents are
discrete.
\end{remark}

\begin{remark}[How the realization reads the quadrant]
\label{rem:conjform}
On every finite divisor-closed coefficient set,
Theorem~\ref{conj:geometric} realizes the successive slope cuts by intervals
in $[0,1]$, the two edges of the quadrant $C_\partial$ by the two endpoints,
and the realizable systems at the two extreme chambers by the extremal systems
$\nu^{\min}$ and $\nu^{\max}$.  In a finite primary
quotient these cuts bound actual open chambers; in the descendent completion
only the algebraic ray factors are asserted to be coefficientwise compatible.
The marginal torus terms of
Theorem~\ref{thm:presentation}(ii) are carried by the $(1,1)$-ray, whose fixed
monomials are exactly its own.  The increasing order of slopes is a choice made in the construction: the same argument realizes the factorization of
$g^{\min\to\max}$ in decreasing order of slopes.
\end{remark}

\subsection{The punctured boundary-degree quadrant}
\label{ssec:missing}

On the punctured quadrant $C_\partial^\times=C_\partial\setminus\{0\}$ only
slope matters.  Items (1), (2), (4) of
Proposition~\ref{thm:positionobstruction7} and the algebraic part of (3) hold
for any finitely many rays through the vertex of a closed strictly convex
cone; what is specific to boundary conditions is that backtracking is realized
by pullback homotopies (Lemma~\ref{lem:polytopeextension}).  For a finite
divisor-closed coefficient set $S$, let $\Dscr_S$ be the union of the wall rays
in the quotient.  A \emph{generic path} in $C_\partial^\times$ has off-wall endpoints and is piecewise smooth and transverse to $\Dscr_S$;
homotopies of generic paths are taken relative endpoints, avoid the vertex,
and pass through finitely many simple tangencies to $\Dscr_S$.  For a compact
polygon $P$, a \emph{$P$-homotopy}, a two-parameter homotopy, is a family of multisections of the
descendent Witten bundles over $P\times\mathcal{PM}_\Gamma$ satisfying the
canonical boundary identities of \cite[Def.~3.22]{GKT} on every parameter
slice, with no transversality required.

\begin{lemma}[Pullback fillers]
\label{lem:polytopeextension}
Let $H=(\mathfrak s_T)_{T\in[T^-,T^+]}$ be a canonical homotopy, let $P$ be a
compact polygon and let $\varphi:P\to[T^-,T^+]$ be smooth.  Then
$(p,x)\mapsto\mathfrak s_{\varphi(p)}(x)$ is a $P$-homotopy, strongly positive
on every parameter slice, whose restriction to $\partial P$ is the
reparametrization of $H$ by $\varphi|_{\partial P}$.  If
$\varphi_0,\varphi_1:P\to[T^-,T^+]$ agree on $\partial P$, the pullbacks along
the linear interpolation from $\varphi_0$ to $\varphi_1$ form a
$(P\times[0,1])$-homotopy joining the two pullbacks relative to $\partial P$.
\end{lemma}

\begin{proof}
The canonical boundary identities and strong positivity are slicewise
conditions, and every slice of a pullback is a slice $\mathfrak s_T$ of $H$; a convex combination of maps to an interval takes
values in it and is constant on $\partial P$, where the maps agree.
\end{proof}

\begin{proposition}[The punctured boundary-degree quadrant]
\label{thm:positionobstruction7}
Let $S$ be finite and divisor-closed.
\begin{enumerate}
\item As a wall-stratified space, $C_\partial^\times$ is isomorphic to
$(0,\infty)\times[0,1]$ with walls the products of $(0,\infty)$ with points,
and it admits a stratum-preserving strong deformation retraction onto a
constant-radius interval.
\item The crossing word of a generic path is determined by its angular
projection.  A homotopy of generic paths inserts or deletes only adjacent inverse
crossings, and every crossing word reduces uniquely to the monotone path in
the chamber-adjacency graph, which is a path.
\item Choose the intermediate families and ray legs supplied by
Theorem~\ref{conj:geometric}.  Assigning $\theta_{\rho,S}$ to a positive
crossing of $\rho$ and $\theta_{\rho,S}^{-1}$ to a negative crossing gives a
reduced automorphism depending only on the angular endpoint chambers.  A
backtracking cancellation is realized by a $P$-homotopy, for a polygon $P$,
pulled back from the corresponding leg, relative to the boundary, and any
two homotopies of this kind are joined by a $(P\times[0,1])$-homotopy of the
same kind.  With $\pi:C_\partial^\times\to\mathbb P(C_\partial)\cong[0,\infty]$ the
projection to the slope space, this reduced path-ordered product factors
through $\pi$: radial paths act by the identity, and positive dilation changes
neither crossing words nor their reduced automorphisms.
\item $C_\partial^\times$ is contractible; every closed generic path has
trivial reduced path-ordered product, and $g^{\min\to\max}$ is the reduced
path-ordered product of a path from one edge of the quadrant to the other.
\end{enumerate}
\end{proposition}

\begin{proof}
(1), (2) and (4).  For $q=q_1e_1+q_2e_2$ put $\varrho(q)=q_1+q_2$ and
$\zeta(q)=q_2/(q_1+q_2)$, an increasing function of the slope.  The map
\begin{equation}
\label{eq:sectorproduct}
 (\varrho,\zeta):C_\partial^\times\xrightarrow{\ \sim\ }(0,\infty)\times[0,1],
 \qquad\text{with inverse}\quad
 (\varrho,\zeta)\longmapsto \varrho\,\bigl((1-\zeta)e_1+\zeta e_2\bigr),
\end{equation}
carries each wall $\rho_i\setminus\{0\}$ to $(0,\infty)\times\{\zeta_i\}$, and
the linear homotopy of $\varrho$ to $1$ with $\zeta$ fixed is the retraction of
(1), so $C_\partial^\times$ is contractible and crossing words are those of the
angular projections.  A homotopy of generic
paths changes a word only at simple tangencies, each inserting or deleting
$\chi_i\chi_i^{-1}$ or $\chi_i^{-1}\chi_i$ ($\chi_i^{\pm1}$ a positive or
negative crossing of the $i$th wall), and crossings of distinct walls cannot
exchange order.  Since the chamber-adjacency graph is a finite path, every
word reduces uniquely to the monotone edge path between its endpoint chambers:
the empty word for a closed path, and the word through all rays, with product
$g^{\min\to\max}$ (Theorem~\ref{thm:canonical}), for a path from the
$\sigma_1$-edge to the $\sigma_2$-edge.

(3) Assign to the $i$th chamber the intermediate family $\mathfrak s_i$ of
Theorem~\ref{conj:geometric} and to $\chi_i^{\pm1}$ the ray leg
$H_i:\mathfrak s_{i-1}\to\mathfrak s_i$, parametrized by $[T_i^-,T_i^+]$, or
its reverse.  The loop $H_i*H_i^{-1}$ is the reparametrization of $H_i$ by a
map $\partial P\to[T_i^-,T_i^+]$ for a polygon $P$; any smooth extension
$\varphi:P\to[T_i^-,T_i^+]$, which exists because the target is an interval,
gives by Lemma~\ref{lem:polytopeextension} a filler
relative to the boundary, and two fillers are joined by the pullback along the
linear interpolation of their maps.  The remaining assertions follow from (2),
since radial paths and positive dilations preserve angular projections.
\end{proof}

\begin{remark}[Interpretation]
\label{rem:interpretation}
The product $g^{\min\to\max}$ from one edge of $C_\partial$ to the other is
not the holonomy of a loop about the vertex: the two edges of $C_\partial$ are
not identified, and on the punctured quadrant the continuation along
homotopies of
boundary conditions sees only slope.  In particular the marginal torus terms
are not identified here with the holonomy of an integral-affine structure.  In
a tropical Landau--Ginzburg model in the sense of \cite{CPS}, the input
includes a singular integral-affine manifold, a polyhedral decomposition and
a multivalued convex piecewise-linear function, and broken lines use that
affine geometry; none of these data is recovered from the quadrant.  Likewise
the present results transport potentials and their decomposition into initial
coefficients, but
they do not construct the operation on pairs of broken lines that would
identify theta-function multiplication constants with open invariants.  Theta
functions in the sense of \cite{GHK15,GHKK,GS-intrinsic} might still span the
Saito--Givental deformation of $x^r+y^s$ and recast the open mirror theorem
\cite[Thm.~0.3]{GKT} as intrinsic mirror symmetry; such a theory would need
data not considered here, for example an energy valuation, a period
normalization, or a logarithmic degeneration with affine gluing.  The
companion paper \cite{annular} constructs, for $x^4+y^4$, a singular
integral-affine annulus over the Carl--Pumperla--Siebert base, with four
singular points, and on it a wall structure consistent at its two joints,
whose singleton walls are the ray factors of $g^{\min\to\max}$ for a single first-descendent marking of twist $(1,0)$ or $(0,1)$, but whose first wall carrying one marking of each of these twists is not; it constructs no theta functions.
\end{remark}

\section{Exact finite calculations}
\label{sec:cert}

The proofs of the structural theorems are independent of software.  Exact
calculations provide the following reproducibility checks; the repository
\url{https://github.com/bwuebben/open-fjrw-scattering} contains the programs
and a file-by-file map from these mathematical statements to their
implementations.

\emph{Critical graphs and boundary rays.}
Direct enumeration from \cite[Not.~2.40]{GKT} agrees with the closed formula
for $N(J,\mathbf d)$, the weight trichotomy, the values of
$N_{\max}(r,s;n)$ and $N_2(r,s)$, and the two-ray example
$J=\{(3,3)^5\}$.  The same calculation gives the ray set
$R(a,b)$ of Proposition~\ref{prop:localfinite}.  Exact symbolic
differentiation verifies bracket closure and the $(1,1)$ exponent shift.

\emph{Periods, closed correlators, and recursion.}
Expansion of the oscillatory integral gives the two-insertion relation
\eqref{eq:seed} and the Gamma functional of Lemma~\ref{lem:periods}.
The closed correlators printed in Example~\ref{ex:x5y5} are obtained
independently from Saito flat coordinates, genus-zero tautological
intersection theory, and Chiodo's formula for Witten's $r$-spin class.  With
these values, both sides of \cite[Thm.~0.5]{GKT} agree in the $72$ and $105$
instances described in Appendix~\ref{sec:x5y5}.

\emph{Ray factors and transport.}
Exact expansion through deformation degree four reproduces the displayed
factors of Theorem~\ref{thm:canonical} in increasing order of slopes,
including the coefficient $-1/750$ of $t_{2,3}t_{3,3}^3X_{1,2}$ on the ray of
$(2,3)$, the coefficient $-1/1250$ of $t_{2,3}^2t_{3,3}^2X_{0,2}$ on the ray
of $(1,3)$, the vanishing coefficient of $t_{3,3}^4$ on the $(1,1)$ ray, and
the value $\nu^{\min}_{(2,7)}=-\tfrac{12}{25}$ of Remark~\ref{rem:tower}.
The transport identities reproduce \eqref{eq:newinvariants}, the initial
coefficient $-25/16$ in Example~\ref{ex:middleseeds}, the foreign
contribution $10/3$ in Example~\ref{ex:farside}, and the bend weights
entering \eqref{eq:seedchain} in that example.

\appendix

\section{Closed FJRW correlators for
\texorpdfstring{$x^{5}+y^{5}$}{x5+y5} and the open topological recursion}
\label{sec:x5y5}

Theorem~\ref{thm:effective} determines the open invariants without closed
input; the closed extended invariants occur as the values of $\Ascr$ on
diagonals without balanced coefficients \cite[Cor.~5.11(3)]{GKT}.  For the
first symmetric hyperbolic Fermat pair $x^{5}+y^{5}$ ($\chat_{W}=\tfrac65>1$)
we compute them exactly.

\begin{example}[Closed vertex weights for $x^{5}+y^{5}$]
\label{ex:x5y5}
By Sebastiani--Thom \cite{FJR13} the state space and cohomological field theory of the Fermat
sum factor as $A_{4}(x)\otimes A_{4}(y)$, the tensor square of the $5$-spin
($A_{4}$) theory; a closed invariant
$\langle\prod_{i}\tau_{d_{i}}^{(a_{i},b_{i})}\rangle$ is the tensor correlator
$\int_{\Mbar_{0,n}}c^{1/5}(\mathbf a)\,c^{1/5}(\mathbf b)\,\psi^{\mathbf d}$ of
the Witten $5$-spin classes $c^{1/5}(\mathbf a)$ and $c^{1/5}(\mathbf b)$ of the
two factors, of degrees $D_{1}=(\sum a_{i}-3)/5$ and
$D_{2}=(\sum b_{i}-3)/5$. We tabulate these with the sign
$(-1)^{D_{1}+D_{2}}$ of the Saito-residue (flat-coordinate)
normalization.  The single-variable $A_{4}$ invariants come from the Saito
primitive form, validated against the $A_{2}$ normalization
$\langle(\tau_{0}^{1})^{4}\rangle=\tfrac13$, WDVV and the $r=2$
Witten--Kontsevich numbers \cite{Kontsevich92}; the correlators with
$D_{1}=0$ or $D_{2}=0$ factorize; and the fully mixed
correlators with $\min(D_{1},D_{2})=1$ are obtained by
reconstructing the class of degree one, a divisor $\sum_{U}c_{U}\delta_{U}$ on
$\Mbar_{0,n}$, from its node-splitting restrictions and pairing it with the
other class, by the genus-zero tautological intersection calculus
\cite{Keel92}.
Representative values (denominators $5^{\,n-3}$; the socle pairing is
$\langle\tau_{0}^{(0,0)}\tau_{0}^{(0,0)}\tau_{0}^{(3,3)}\rangle=1$):
\[
\renewcommand{\arraystretch}{1.4}
\begin{array}{lll}
\langle(\tau_{0}^{(2,1)})^{3}\tau_{0}^{(2,0)}\rangle=-\tfrac25, &
\langle\tau_{0}^{(1,3)}\tau_{0}^{(1,0)}(\tau_{0}^{(3,0)})^{2}\rangle=-\tfrac15
& (D_{1},D_{2})=(1,0),\\
\langle(\tau_{0}^{(2,1)})^{2}\tau_{0}^{(2,2)}(\tau_{0}^{(1,2)})^{2}\rangle
=\tfrac{4}{25}, &
\langle(\tau_{0}^{(2,2)})^{3}(\tau_{0}^{(1,1)})^{2}\rangle=\tfrac{2}{25}
& (1,1),\\
\langle(\tau_{0}^{(2,3)})^{2}\tau_{0}^{(1,3)}\tau_{0}^{(1,2)}
(\tau_{0}^{(1,1)})^{2}\rangle=-\tfrac{4}{125}, &
\langle(\tau_{0}^{(2,3)})^{2}(\tau_{0}^{(1,3)})^{3}\tau_{0}^{(1,2)}
\tau_{0}^{(0,1)}\rangle=\tfrac{6}{625}
& (1,2),\ (1,3).
\end{array}
\]
In the smallest case with both classes of codimension two, $n=7$, the
codimension-two stratum intersection form gives for instance
$\langle\tau_{0}^{(3,1)}\tau_{0}^{(2,1)}(\tau_{0}^{(2,2)})^{3}\tau_{0}^{(1,2)}
\tau_{0}^{(1,3)}\rangle=\tfrac{34}{625}$.  The
closed extended invariants in the normalization of \cite[Thm.~0.3]{GKT} are
obtained by removing the sign $(-1)^{D_1+D_2}$, for instance
$\langle(\tau_{0}^{(2,1)})^{3}\tau_{0}^{(2,0)}\rangle^{\mathrm{ext}}=\tfrac25$.
\end{example}

Substituting these closed values into the open topological
recursion of \cite[Thm.~0.5]{GKT}, with the values of $\Ascr$ given by
\cite[Cor.~5.11]{GKT}, gives equality in the $72$ instances with two internal
markings of twists in $\{(1,1),(2,1),(1,2),(2,2),(1,3),(3,1),(2,3),(3,2)\}$
and in the $105$ instances with three internal markings of twists in
$\{(1,1),(2,1),(1,2),(2,2),(1,3)\}$, all without descendents and including the
extended sectors, in which a twist component equals $-1$ (a Ramond twist).  The recursion is invariant under
the sign change above; it is a consistency check of the closed correlators and
involves no open invariant.  An evaluation of Witten's $r$-spin class through
Chiodo's formula \cite{Chiodo08} independently reproduces the values in
Example~\ref{ex:x5y5}, including $\tfrac{34}{625}$.

\begin{thebibliography}{GKT26}

\bibitem[BCT24]{BCT21} A.~Buryak, E.~Clader, R.~J.~Tessler,
\emph{Open $r$-spin theory II: The analogue of Witten's conjecture for
$r$-spin disks}, J.\ Differential Geom.\ 128 (2024), no.~1, 1--75;
arXiv:1809.02536.

\bibitem[Bri17]{Bridgeland} T.~Bridgeland,
\emph{Scattering diagrams, Hall algebras and stability conditions},
Algebr.\ Geom.\ 4 (2017), no.~5, 523--561; arXiv:1603.00416.

\bibitem[Chi08]{Chiodo08} A.~Chiodo,
\emph{Towards an enumerative geometry of the moduli space of twisted curves
and $r$th roots}, Compos.\ Math.\ 144 (2008), no.~6, 1461--1496.

\bibitem[CPS24]{CPS} M.~Carl, M.~Pumperla, B.~Siebert,
\emph{A tropical view on Landau--Ginzburg models},
Acta Math.\ Sin.\ (Engl.\ Ser.)\ 40 (2024), no.~1, 329--382;
arXiv:2205.07753.

\bibitem[FJR13]{FJR13} H.~Fan, T.~Jarvis, Y.~Ruan,
\emph{The Witten equation, mirror symmetry, and quantum singularity theory},
Ann.\ of Math.\ 178 (2013), no.~1, 1--106; arXiv:0712.4021.

\bibitem[GHK15]{GHK15} M.~Gross, P.~Hacking, S.~Keel,
\emph{Mirror symmetry for log Calabi--Yau surfaces I},
Publ.\ Math.\ Inst.\ Hautes \'Etudes Sci.\ 122 (2015), 65--168;
arXiv:1106.4977.

\bibitem[GHKK18]{GHKK} M.~Gross, P.~Hacking, S.~Keel, M.~Kontsevich,
\emph{Canonical bases for cluster algebras}, J.\ Amer.\ Math.\ Soc.\ 31
(2018), no.~2, 497--608; arXiv:1411.1394.

\bibitem[GKT22]{GKT} M.~Gross, T.~L.~Kelly, R.~J.~Tessler,
\emph{Open FJRW theory and mirror symmetry}, Geom.\ Topol.\ 30 (2026),
no.~8, 2779--2960; doi:10.2140/gt.2026.30.2779;
arXiv:2203.02435v3 (2022; revised 2026).

\bibitem[GKT26]{GKTsurvey} M.~Gross, T.~L.~Kelly, R.~J.~Tessler,
\emph{Open enumerative geometries for Landau--Ginzburg models},
arXiv:2602.12707v1 (2026).

\bibitem[GPS10]{GPS10} M.~Gross, R.~Pandharipande, B.~Siebert,
\emph{The tropical vertex}, Duke Math.\ J.\ 153 (2010), no.~2, 297--362;
arXiv:0902.0779.

\bibitem[Gro10]{Gross10} M.~Gross,
\emph{Mirror symmetry for $\PP^{2}$ and tropical geometry},
Adv.\ Math.\ 224 (2010), no.~1, 169--245; arXiv:0903.1378.

\bibitem[GS22]{GS-intrinsic} M.~Gross, B.~Siebert,
\emph{The canonical wall structure and intrinsic mirror symmetry},
Invent.\ Math.\ 229 (2022), no.~3, 1101--1202; arXiv:2105.02502.

\bibitem[HLSW22]{HLSW} W.~He, S.~Li, Y.~Shen, R.~Webb,
\emph{Landau--Ginzburg mirror symmetry conjecture},
J.\ Eur.\ Math.\ Soc.\ 24 (2022), no.~8, 2915--2978; arXiv:1503.01757.

\bibitem[Kee92]{Keel92} S.~Keel,
\emph{Intersection theory of moduli space of stable $n$-pointed curves of
genus zero}, Trans.\ Amer.\ Math.\ Soc.\ 330 (1992), no.~2, 545--574.

\bibitem[Kon92]{Kontsevich92} M.~Kontsevich,
\emph{Intersection theory on the moduli space of curves and the matrix Airy
function}, Comm.\ Math.\ Phys.\ 147 (1992), no.~1, 1--23.

\bibitem[KS06]{KS06} M.~Kontsevich, Y.~Soibelman,
\emph{Affine structures and non-Archimedean analytic spaces}, in
\emph{The unity of mathematics}, Progr.\ Math.\ 244, Birkh\"auser (2006), 321--385;
arXiv:math/0406564.

\bibitem[LLSS17]{LLSS} C.~Li, S.~Li, K.~Saito, Y.~Shen,
\emph{Mirror symmetry for exceptional unimodular singularities},
J.\ Eur.\ Math.\ Soc.\ 19 (2017), no.~4, 1189--1229; arXiv:1405.4530.

\bibitem[Mah24]{Maher} R.~Maher,
\emph{Predictions in open Fan--Jarvis--Ruan--Witten theory via mirror symmetry,
modularity and wall-crossing}, Ph.D.\ thesis, University of Birmingham, 2024;
\texttt{etheses.bham.ac.uk/id/eprint/15556}.

\enlargethispage{2\baselineskip}
\bibitem[Wue26a]{note} B.~J.~Wuebben,
\emph{The central-charge threshold for wall-crossing in two-variable open FJRW
theory}, companion note (2026),
\url{https://github.com/bwuebben/open-fjrw-scattering}.

\bibitem[Wue26b]{annular} B.~J.~Wuebben,
\emph{Annular scattering in quartic open FJRW theory}, manuscript (2026),
\url{https://github.com/bwuebben/open-fjrw-scattering}.

\end{thebibliography}

\bigskip
\noindent{\sc Bernd J. Wuebben, New York, NY,} \texttt{wuebben@gmail.com}

\end{document}
